\chapter{The standard solution}
\label{chap:standard-solution}
\label{stdsolnsect}

The process of surgery involves making a choice of the metric on a
three-ball to `glue in'. In order to match approximatively with the
metric coming from the flow, the metric we glue in must be
asymptotic to the product of a round two-sphere and an interval near
the boundary. There is no natural choice for this metric; yet it is
crucial to the argument that we choose an initial metric so that the
Ricci flow with these initial conditions has several properties.
Conditions on the initial metric that  ensure the required
properties for the subsequence flow are contained in the following
definition.


\begin{definition}[Standard initial metric]
\label{def:standard-initial-metric}
\label{defn:stinmetr}
\uses{def:sectional-curvature}
A {\em standard initial metric}
is a metric $g_0$ on $\Ar^3$ with the following properties:
\begin{itemize}
\item $g_0$ is a complete metric.
\item $g_0$ has non-negative sectional curvature at every point.
\item $g_0$ is invariant under the usual $SO(3)$-action on $\Ar^3$.
\item there is a compact ball $B\subset \Ar^3$ so
that the restriction of the metric $g_0$ to the complement of this
ball is isometric to the product $(S^2,h)\times (\Ar^+,ds^2)$ where
$h$ is the round metric of scalar curvature $1$ on $S^2$.
\item $g_0$ has constant sectional curvature $1/4$ near the origin. (This point will be
denoted $p$ and is called the {\em tip} of the initial metric.)
\end{itemize}
\end{definition}

Actually, one can work with an alternative weaker version of the
fourth condition, namely:

\noindent (iv) $g_0$ is asymptotic at infinity in the
$C^\infty$-topology to the product of the round metric $h_0$ on
$S^2$ of  scalar curvature $1$ with the usual metric $ds^2$ on the
real line. By this we mean  that if $x_n\in \Ar^3$ is any sequence
converging to infinity, then the based Riemannian manifolds
$(\Ar^3,g_0,x_n)$ converge smoothly to $(S^2,h_0)\times (\Ar,ds^2)$.
But we shall only use standard initial metrics as given in
Definition~\ref{def:standard-initial-metric}.




\begin{lemma}[Existence of a standard initial metric]
\label{lem:existence-standard-initial-metric}
\label{stdinit}
\uses{def:standard-initial-metric}
There is a standard initial metric.
\end{lemma}

\begin{proof}
\uses{def:standard-initial-metric}
We construct our Riemannian manifold as follows. Let
$(x_0,x_1,x_2,x_3)$ be Euclidean coordinates on $\Ar^4$. Let
$y=f(s)$ be a function defined for $s\ge 0$ and satisfying:
\begin{enumerate}
\item[(1)] $f$ is $C^\infty$ on $(0,\infty)$
\item[(2)] $f(s)>0$ for all $s>0$.
\item[(3)] $f''(s)\le 0$ for all $s>0$.
\item[(4)] There is $s_1>0$ such that $f(s)=\sqrt{2}$ for all $s\ge
s_1$.
\item[(5)] There is $s_0>0$ such that
$f(s)=\sqrt{4s-s^2}$ for all $s\in [0,s_0]$.
\end{enumerate}
Given such a function $f$, consider the graph
$$\Gamma=\{(x_0,x_1)\bigl|\bigr. x_0\ge 0 \text{ and }
x_1=f(x_0)\}$$ in the $(x_0,x_1)$-plane. We define $\Sigma(f)$ by
rotating $\Gamma$ about the $x_0$-axis in four-space:
$$\Sigma(f)=\{(x_0,x_1,x_2,x_3)\bigl|\bigr. x_0\ge 0 \text{ and }
x_1^2+x_2^2+x_3^2=f(x_0)^2\}.$$ Because of the last condition on
$f$, there is a neighborhood of $0\in \Sigma(f)$ that is isometric
to a neighborhood of the north pole in the three-sphere of radius
$2$. Because of this and the first item, we see that $\Sigma(f)$ is
a smooth submanifold of $\Ar^4$. Hence, it inherits a Riemannian
metric $g_0$. Because of the fourth item, a neighborhood of infinity
of $(\Sigma(f),g_0)$ is isometric to $(S^2,h)\times (0,\infty)$, and
in particular, $(\Sigma(f),g_0)$ is complete. Clearly, the rotation
action of $S0(3)$ on $\Sigma(f)$, induced by the orthogonal action
on the last three coordinates in $\Ar^4$, is an isometric action
with the origin as the only fixed point. It is also clear that
$\Sigma(f)$ is diffeomorphic to $\Ar^3$ by a diffeomorphism that
send the $SO(3)$ action to the standard one on $\Ar ^3$.

It remains to compute the sectional curvatures of $g_0$. Let $q\in
\Sigma(f)$ be a point distinct from the fixed point of the
$SO(3)$-action. Direct computation shows that the tangent plane to
the two-dimensional $SO(3)$-orbit through $q$ is a principal
direction for the curvature, and the sectional curvature on this
 tangent two-plane  is given by
$$\frac{1}{f(q)^2(1+f'(q)^2)}.$$
On the subspace in $\wedge^2T_q\Sigma(f)$ perpendicular to the line
given by this two-plane, the curvature is constant  with eigenvalue
$$\frac{-f''(q)}{f(q)(1+f'(q)^2)^2}.$$
Under our assumptions about $f$, it is clear that $\Sigma(f)$ has
non-negative curvature and has constant sectional curvature $1/4$
near the origin. It remains to choose the function $f$ satisfying
Items (1) -- (5) above.


Consider the function $h(s)=(2-s)/\sqrt{4s-s^2}$. This function is
integrable from $0$ and the definite integral from  zero to $s$ is
equal to $\sqrt{4s-s^2}$. Let $\lambda(s)$ be a non-increasing
$C^\infty$-function defined on $[0,1/2]$, with $\lambda$ identically
one near $0$ and identically equal to $0$ near $1/2$. We extend
$\lambda$ to be identically $1$ for $s<0$ and identically $0$ for
$s>1/2$. Clearly,
$$\int_0^2h(s)\lambda(s-3/2)ds>\int_0^{3/2}h(s)ds>\sqrt{2}$$
and $$\int_0^2h(s)\lambda(s)ds<\int_0^{1/2}h(s)<\sqrt{2}.$$ Hence,
for some $s_0\in (1/2,3/2)$ we have
$$\int_0^2h(s)\lambda(s-s_0)ds=\sqrt{2}.$$
We define $$f(s)=\int_0^sh(\sigma)\lambda(\sigma-s_0)d\sigma.$$ It
is easy to see that $f$ satisfies all the above conditions.
\end{proof}

The following lemma is clear from the construction.

\begin{lemma}[Asymptotic and curvature bounds for the standard initial metric]
\label{lem:standard-initial-metric-asymptotic-bounds}
\label{A_0}
\uses{lem:existence-standard-initial-metric}
There is $A_0<\infty$ such that
$$(\Ar^3\setminus B(0,A_0),g(0))$$
is isometric to the product of a round metric on $S^2$ of scalar
curvature $1$ with the Euclidean metric on $[0,\infty)$. There is a
constant $K<\infty$ such that the volume of $B_{g(0)}(0,A_0)$ is at
most $K$. Furthermore, there is a constant $D<\infty$ so that the
scalar curvature of standard initial metric $(\Ar^3,g(0))$ is
bounded above by $D$ and below by $D^{-1}$.
\end{lemma}



\subsection{Uniqueness and properties: The statement}

{\bf Fix once and for all a standard initial metric $g_0$ on
$\Ar^3$.}

\begin{definition}[Partial and standard Ricci flow]
\label{def:standard-ricci-flow}
\uses{def:standard-initial-metric,def:ricci-flow}
A {\em partial standard Ricci flow} is a Ricci flow $(\Ar^3,g(t)),\ 0\le t<T$, such
that $g(0)=g_0$ and such that the curvature is locally bounded in
time. We say that a partial standard Ricci flow is a {\em standard
Ricci flow} if it has the property
that $T$ is maximal in the sense that there is no extension of the
flow to a flow on a larger time interval $[0,T')$ with $T'>T$ with
the property that the extension has curvature locally bounded in
time.
\end{definition}

Here is the main result of this chapter.

\begin{theorem}[Existence, uniqueness, and properties of the standard solution]
\label{thm:standard-solution-existence-uniqueness}
\label{stdsoln}
\uses{def:standard-ricci-flow,def:epsilon-neck}
There is a standard Ricci flow defined for some
positive amount of time. Let $(\Ar^3,g(t)),\ 0\le t<T$, be a
standard Ricci flow. Then the following hold.
\begin{enumerate}
\item[(1)] {\bf (Uniqueness):}
If $(\Ar^3,g'(t)),\ 0\le t<T'$, is a standard Ricci flow, then
$T'=T$ and $g'(t)=g(t)$.
\item[(2)] {\bf (Time Interval):} $T=1$.
\item[(3)] {\bf (Positive curvature):} For each $t\in (0,1)$ the metric $g(t)$ on $\Ar^3$ is complete
of strictly positive curvature.
\item[(4)] {\bf ($SO(3)$-invariance):} For each $t\in [0,T)$ the
Riemannian manifold $(\Ar^3,g(t))$  is invariant under the
$SO(3)$-action on $\Ar^3$.
\item[(5)] {\bf (Asymptotics at $\infty$):}
For any $t_0<1$ and any $\epsilon>0$ there is a compact subset $X$
of $\Ar^3$ such that for any $x\in \Ar ^3\setminus X$ the
restriction of the standard flow to an appropriate neighborhood of
$x$ for time $t\in [0,t_0]$ is within $\epsilon$ in the
$C^{[1/\epsilon]}$-topology of the product Ricci flow $(S^2\times
(-\epsilon^{-1},\epsilon^{-1})),h(t)\times ds^2,\ 0\le t\le t_0$,
where $h(t)$ is the round metric with scalar curvature $1/(1-t)$ on
$S^2$.
\item[(6)] {\bf (Non-collapsing):} There are $r>0$ and $\kappa>0$ such that $(\Ar^3,g(t)),\ 0\le
t<1$, is $\kappa$-non-collapsed on scales less than $r$.
\end{enumerate}
\end{theorem}

The proof of this result occupies the next few subsections. All the
properties except the uniqueness are fairly straightforward to
prove. We establish uniqueness by reducing the Ricci flow to the
Ricci-DeTurck by establishing the existence of a solution to the
harmonic map flow in this case. This technique can be made to work
more generally in the case of complete manifolds of bounded
curvature, see \cite{ChenZhu}, but we preferred to give the more
elementary argument that covers this case, where the symmetries
allow us to reduce the existence problem for the harmonic map flow
to a problem that is the essentially one-dimensional. Also, in the
rest of the argument one does not need uniqueness, only a
compactness result for the space of all Ricci flows of bounded
curvature on each time-slice with the given initial conditions.
Kleiner and Lott pointed out to us that this uniqueness can be
easily derived from the other properties by arguments similar to
those used to establish the compactness of the space of
$\kappa$-solutions.

\section{Existence of a standard flow}

For any $R<\infty$, denote by $B_R\subset \Ar^3$, the ball of radius
$R$ about the origin in the metric $g_0$. For $R\ge A_0+1$,  a
neighborhood of the boundary of this ball is isometric to
$(S^2,h)\times ([0,1],ds^2)$. Thus, in this case, we can double the
ball, gluing the boundary to itself by the identity, forming a
manifold we denote by $S^3_R$. The doubled metric will be a smooth
Riemannian metric $g_R$ on $S^3_R$. Let $p\in S^3_R$ be the image of
the origin in the first copy of $B_R$. Now take a sequence, $R_n$,
tending to infinity to construct based Riemannian manifolds
$(S^3_{R_n},g_{R_n},p)$ that converge geometrically to
$(\Ar^3,g_0,p)$. For each $n$, let $(S^3_{R_n},g_{R_n}(t)),\ 0\le
t<T_n$ be maximal Ricci flow with $(S^3_{R_n},g_{R_n})$ as initial
metric.
 The maximum principle
applied to Equation~(\ref{Revol}), $\partial R/\partial t=\triangle
R+|\Ric|^2$, then implies by Proposition~\ref{prop:forward-difference-maximum} that
the maximum of $R$ at time $t$, $R_{\max}(t)$ obeys the
inequality $\partial R_{\max}/\partial t\le R_{\max}(t)^2$,
and integrating this inequality (i.e., invoking
Lemma~\ref{lem:forward-difference-quotient}) one finds a positive constants $t_0$ and
$Q_0$ such that for each $n$, the norm of the scalar curvature of
$g_{R_n}(t)$ are bounded by $Q_0$ on the interval $[0,\max(t_0,T_n))$. By Corollary~\ref{rm>0}, for each $n$ the sectional
curvature of the flow $(S^3_{R_n},g_{R_n}(t)),\ 0\le t<T_n$ is
non-negative, and hence the sectional curvature of  this flow is
also bounded by $Q_0$ on $[0,\max(t_0,T_n))$. It now follows
from Proposition~\ref{flowextend} and the fact that the $T_n$ are
maximal that $T_n>t_0$ for all $n$.
 Since the Riemann
curvatures of the $(S^3_{R_n},g_{R_n}(t)),\ 0\le t<t_0$, are bounded
independent of $n$, and since the $(S^3_{R_n},g_{R_n},p)$ converge
geometrically to $(\Ar^3,g_0,p)$, it follows from
Theorem~\ref{flowlimit} that there is a geometric limiting flow
defined on $[0,t_0)$. Since this flow is the geometric limit of
flows of uniformly bounded curvature, it has uniformly bounded
curvature. Taking a maximal extension of this flow to one of locally
bounded curvature gives a standard flow.



\section{Completeness, positive curvature, and asymptotic behavior}

 Let $\left( \mathbb{R}^{3},g(t)
\right)$, $t\in [0,T)$, be a partial standard solution. Let
$y_{i}\rightarrow\infty$ be the sequence of points in
$\mathbb{R}^{3}$ converging to infinity. From the  definition we see
that the based Riemannian manifolds $\left(
\mathbb{R}^{3},g_{0},y_{i}\right)$ converge smoothly to
$(S^2\times\mathbb{R},h(0)\times ds^2)$ where $h(0)$ is the round
metric of scalar curvature $1$ on $S^2$.

Let us begin by proving the third item in the statement of
Theorem~\ref{thm:standard-solution-existence-uniqueness}:

\begin{lemma}[Completeness and positive curvature of the standard flow]
\label{lem:standard-flow-complete-positive-curvature}
\label{third}
\uses{def:standard-ricci-flow}
For each $t_0\in [0,T)$ the Riemannian manifold $(\Ar^3,g(t_0))$ is
complete and of positive curvature.
\end{lemma}


\begin{proof}
\uses{prop:forward-difference-maximum,lem:curvature-evolution-frame,nullspace}
Fix $t_0\in [0,T)$. By hypothesis $(\Ar^3,g(t)),\ 0\le t\le t_0$ has bounded
curvature. Hence, there is a constant $C<\infty$ such that $ g(0)\le Cg(t_0)$,
so that for any points $x,y\in \Ar^3$, we have $d_0(x,y)\le
\sqrt{C}d_{t_0}(x,y)$. Since $g(0)$ is complete, this implies that $g(t_0)$ is
also complete.

Now let us show that $(M,g(t_0))$ has non-negative curvature. Here,
the argument is the analogue of the proof of Corollary~\ref{lem:curvature-evolution-frame}
with one additional step, the use of a function $\varphi$ to
localize the argument. Suppose this is false, i.e., suppose that
there is $x\in M$ with $\Rm(x,t_0)$ having an eigenvalue less
than zero. Since the restriction of the flow to $[0,t_0]$ is
complete and of bounded curvature, according to
\cite{Hamiltonlimits} for any constants $C<\infty$ and $\eta>0$ and
any compact subset $K\subset M\times [0,t_0]$ there is $\epsilon>0$
and a function $\varphi\colon M\times [0,t_0]\to \Ar$ with the
following properties:
%%%%
\begin{enumerate}
\item[(1)] $\varphi|_K\le \eta$.
\item[(2)] $\varphi\ge \epsilon$ everywhere.
\item[(3)] For each $t\in [0,t_0]$ the restriction of $\varphi$  to $M\times \{t\}$
goes to infinity at infinity in the sense that  for any $A<\infty$
the pre-image $\varphi^{-1}([0,A]\cap (M\times\{t\})$ is compact.
\item[(4)] On all of $M\times [0,t_0]$ we have
$
  \left(\frac{\partial}{\partial t}-\triangle\right)\varphi\ge C\varphi.
$
\end{enumerate}
%%%%
Recall from Section~\ref{4.3.3} that ${\mathcal T}$ is the curvature
tensor written with respect to an evolving orthonormal frame
$\{F_\alpha\}$ for the tangent bundle. Consider the symmetric,
horizontal two-tensor $\widehat{\mathcal T}={\mathcal T}+\varphi g$.
Let $\hat\mu(x,t)$ denote the smallest eigenvalue of this symmetric
two-tensor at $(x,t)$. Clearly, since the curvature is bounded, it
follows from the third property of $\varphi$ that for each $t\in
[0,t_0]$ the restriction of $\hat\mu$ to $M\times\{t\}$ goes to
infinity at infinity in $M$. In particular, the subset of $(x,t)\in
M\times [0,t_0]$ with the property that $\hat\mu(x,t)\le
\hat\mu(y,t)$ for all $y\in M$ is a compact subset of $M\times
[0,t_0]$. It follows from Proposition~\ref{prop:forward-difference-maximum} that
$f(t)=\min_{x\in M}\hat\mu(x,t)$ is a continuous function of
$t$. Choosing $\eta>0$ sufficiently small and $K$ to include
$(x,t_0)$, then $\widehat{\mathcal T}$ will have a negative
eigenvalue at $(x,t_0)$. Clearly, it has only positive eigenvalues
on $M\times \{0\}$. Thus, there is $0<t_1<t_0$ so that
$\widehat{\mathcal T}$ has only positive eigenvalues on $M\times
[0,t_1)$ but has a zero eigenvalue at $(y,t_1)$ for some $y\in M$.
That is to say, ${\mathcal T}\ge -\varphi g$ on $M\times [0,t_1]$.
Diagonalizing ${\mathcal T}$ at any point $(x,t)$ with $t\le t_1$,
all its eigenvalues are at least $-\varphi(x,t_1)$. It follows
immediately that on $M\times [0,t_1]$ the smallest eigenvalue of the
symmetric form ${\mathcal T}^2+{\mathcal T}^\#$ is bounded below by
$2\varphi$. Thus, choosing $C\ge 4$  we see that for $t\le t_1$
every eigenvalue of ${\mathcal T}^2+{\mathcal T}^\#$ is at least
$-C\varphi/2$.


We compute the evolution equation using the formula in
Lemma~\ref{lem:curvature-evolution-frame} for the evolution of ${\mathcal T}$ in an evolving
orthonormal frame:
\begin{eqnarray*}\frac{\partial \widehat{\mathcal T}}{\partial t} & = &
\frac{\partial{\mathcal T}}{\partial t}+\frac{\partial
\varphi}{\partial t}g-2\varphi \Ric(g) \\
& = & \triangle {\mathcal T}+{\mathcal T}^2+{\mathcal T}^\#+\frac{\partial \varphi}{\partial t}g-2\varphi \Ric(g) \\
& = & \triangle \widehat{\mathcal T}+{\mathcal T}^2+{\mathcal
T}^\#+\left(\frac{\partial \varphi}{\partial t}-\triangle
\varphi\right)g-2\varphi \Ric(g)\\
&\ge & \triangle \widehat{\mathcal T}+{\mathcal T}^2+{\mathcal
T}^\#+\left(Cg-2\Ric(g)\right)\varphi.
\end{eqnarray*}
Since every eigenvalue of ${\mathcal T}^2+{\mathcal T}^\#$  on
$M\times [0,t_1]$ is at least $-C\varphi/2$, it follows that on
$M\times [0,t_1]$
$$\frac{\partial \widehat{\mathcal T}}{\partial t}\ge
 \triangle \widehat{\mathcal T}+(Cg/2-2\Ric_g)\varphi.$$ Once again assuming that $C$ is
sufficiently large, we see that for any $t\le t_1$
$$\frac{\partial \widehat{\mathcal T}}{\partial t}\ge \triangle
\widehat {\mathcal T}.$$ Thus, at any local minimum $x\in M$ for
$\hat\mu(\cdot,t)$, we have
$$\frac{\partial \hat\mu}{\partial t}\ge 0.$$
This immediately implies by Proposition~\ref{prop:forward-difference-maximum} that
$\psi(t)=\min_{x\in M}\hat\mu(x,t)$ is a non-decreasing
function of $t$. Since its value at $t=0$ is at least $\epsilon>0$
and its value at $t_1$ is zero, this is a contradiction. This
establishes that the solution has non-negative curvature everywhere.
Indeed, by Corollary~\ref{nullspace} it has strictly positive
curvature for every $t>0$.
\end{proof}





Now let us turn to the asymptotic behaviour of the flow.


Fix $T'<T$. Let $y_k$ be a sequence tending to infinity in
$(\Ar^3,g_0)$. Fix $R<\infty$. Then there is $k_0(R)$ such that for
all $k\ge k_0(R)$ there is an isometric embedding $\psi_k\colon
(S^2,h)\times (-R,R)\to (\Ar^3,g_0)$ sending $(x,0)$ to $y_k$. These
maps realize the product $(S^2,h)\times (\Ar,ds^2)$ as the geometric
limit of the $(\Ar^3,g_0,y_i)$. Furthermore, for each $R<\infty$
there is a uniform $C^\infty$ point-wise bound to the curvatures of
$g_0$ restricted to the images of the $\psi_k$ for $k\ge k_0(R)$.
 Since the flow $g(t)$ has bounded
curvature on $\Ar^3\times [0,T']$, it follows from
Theorem~\ref{shiw/deriv} that there are uniform $C^\infty$
point-wise bounds for the curvatures of $g(t)$ restricted to
$\psi_k(S^2\times (-R,R))$. Thus,
 by Theorem~\ref{partialflowlimit}, after passing to a subsequence, the flows
$\psi_k^*g(t)$ converge to a  limiting flow on $S^2\times \Ar$. Of
course, since the curvature of $g(t)$ is everywhere $\ge 0$, the
same is true of this limiting flow. Since the time-slices of this
flow have two ends, it follows from Theorem~\ref{thm:splitting-theorem} that
every manifold in the flow is a product of a compact surface with
$\Ar$. According to Corollary~\ref{nullspace} this implies that the
flow is the product $(S^2,h(t))\times (\Ar,ds^2)$. This means that
given $\epsilon>0$, for all $k$ sufficiently large, the restriction
of the flow to the cylinder of length $2R$ centered at $y_k$  is
within $\epsilon$ in the $C^{[1/\epsilon]}$-topology of the
shrinking cylindrical flow on time $[0,T']$. Given $\epsilon>0$ and
$R<\infty$ this statement is true for all $y$ outside a compact ball
$B$ centered at the origin.

We have now established the following

\begin{proposition}[Asymptotic convergence to the standard evolving cylinder]
\label{prop:standard-flow-asymptotic-cylinder}
\label{item5}
\uses{def:standard-ricci-flow,ex:shrinking-round-cylinder}
Given $T'<T$ and $\epsilon>0$ there is a compact ball $B$ centered
at the origin of $\Ar^3$ such that the restriction of the flow
$(\Ar^3\setminus B,g(t)),\ 0\le t\le T'$, is within $\epsilon$ in
the $C^{[1/\epsilon]}$-topology of the standard evolving cylinder
$(S^2,h(t))\times (\Ar^+,ds^2)$.
\end{proposition}

\begin{corollary}[Maximal time is at most one]
\label{cor:standard-flow-time-bound-one}
\label{1}
\uses{prop:standard-flow-asymptotic-cylinder}
The maximal time $T$ is $\le 1$.
\end{corollary}

\begin{proof}
\uses{prop:standard-flow-asymptotic-cylinder}
If $T>1$, then we can apply the above result to $T'$ with
$1<T'<T$,and see that the solution at infinity is asymptotic to the
evolving cylinder $(S^2,h(t))\times (\Ar,ds^2)$ on the time interval
$[0,T']$. But this is absurd since this evolving cylindrical flow
becomes completely singular at time $T=1$.
\end{proof}







\section{Standard solutions are rotationally symmetric}

Next, we consider the fourth item in the statement of the theorem.
Of course, rotational symmetry would follow immediately from
uniqueness. But here we shall use the rotational symmetry to reduce
the uniqueness problem to a one-dimensional problem which we then
solve. One can also use the general uniqueness theorem for complete,
non-compact manifolds due to Chen and Zhu (\cite{ChenZhu}), but we
have chosen to present a more elementary, self-contained argument in
this special case which we hope will be more accessible.

Let $\Ric_{ij}$ be the Ricci tensor and $\Ric^i_k= g^{ij}
\Ric_{jk}$ be the dual tensor. Let $X$ be a vector field
evolving by
\begin{equation}\label{eq1}\frac{\partial}{\partial t}X=\triangle X+\Ric(X,\cdot)^*.
\end{equation}
 In local coordinates $(x^1,\ldots,x^n)$, if
 $X=X^i\partial_i$, then the equation becomes
\begin{equation}
\frac{\partial}{\partial t}X^{i}=(\Delta X)^{i}+\Ric_{k}^{i}X^{k} \label{eq.3}
\end{equation}
Let $X^*$ denote the dual one-form to $X$. In local coordinates we
have $X^*=X^*_idx^i$ with $X^*_i=g_{ij}X^j$. Since the evolution
equation for the metric is the Ricci flow, the evolution equation
for $X^*$ is
$$\frac{\partial X^*}{\partial t}=\triangle X^*-\Ric(X,\cdot),$$
or in local coordinates
$$\frac{\partial X^*_i}{\partial t}=(\triangle X^*)_i-\Ric_{ij}X^j.$$

\begin{lemma}[Evolution of $V=\nabla X^*$]
\label{lem:v-tensor-evolution}
\label{12.9}
\uses{def:standard-ricci-flow}
With $X$ and its dual $X^*$ evolving by the above equations, set
$V=\nabla X^*$, so that $V$ is a contravariant two-tensor. In local
coordinates we have $V=V_{ij}dx^i\otimes dx^j$ with
$$V_{ij}=(\nabla_iX)_j=g_{jk}(\nabla_iX)^k.$$
This symmetric two-tensor satisfies
\begin{equation}
\frac{\partial}{\partial t}V=\triangle
V-\left(2{{R_k}^{rl}}_jV_{rl}+\Ric_k^lV_{lj} +\Ric_j^lV_{kl}\right)dx^k\otimes dx^j. \label{eq.4}
\end{equation}
\end{lemma}

\begin{remark}[Covariant derivative and curvature acting on one-forms]
\label{rem:covariant-derivative-one-forms}
\uses{def:riemann-curvature-tensor}
The covariant derivative acts on one-forms $\omega$ in such a way
that the following equation holds:
$$\langle\nabla(\omega),\xi\rangle=\langle
\omega,\nabla(\xi)\rangle$$ for every vector field $\xi$. This means
that in local coordinates we have
$$\nabla_{\partial_r}(dx^k)=-\Gamma_{rl}^kdx^l.$$ Similarly, the Riemann
curvature acts on one-forms $\omega$ satisfying
$$\Rm(\xi_1,\xi_2)(\omega)(\xi)=-\omega\left(\Rm(\xi_1,\xi_2)(\xi)\right).$$
Recall that in local coordinates  $$R_{ijkl}=\langle \Rm(\partial_i,\partial_j)(\partial_l),\partial_k\rangle.$$ Thus, we
have
$$\Rm(\partial_i,\partial_j)(dx^k)=-g^{ka}R_{ijal}dx^l=-{{R_{ij}}^k}_ldx^l,$$
where as usual we use the inverse metric tensor to raise the index.


Also, notice that  $\Delta X_{i}-\Ric_{ik}X^{k}=-\Delta_{d}X_{i\text{ }}$, where by $\Delta_d$ we
mean the Laplacian associated to the operator $d$ from vector fields
to one-forms with values in the vector field. Since
\begin{align*}
 -\left(  d\delta+\delta d\right)  X_{i}
&  =-\nabla_{i}\left(  -\nabla^{k}X_{k}\right)  -\left(
-\nabla^{k}\right)
\left(  \nabla_{k}X_{i}-\nabla_{i}X_{k}\right) \\
&
=\nabla_{i}\nabla^{k}X_{k}+\nabla^{k}\nabla_{k}X_{i}-\nabla^{k}\nabla
_{i}X_{k}\\
&  ={{{R_{i}}^k}_k}^jX_{j}+\nabla^{k}\nabla_{k}X_{i}=\Delta
X_{i}-\Ric_i^jX_{j}.
\end{align*}
\end{remark}

\begin{proof}(of Lemma~\ref{lem:v-tensor-evolution})
The computation is routine, if complicated. We make the computation
at a point $(p,t)$ of space-time. We fix local $g(t)$-Gaussian
coordinates $(x^1,\ldots,x^n)$ centered at $p$ for space, so that
the Christoffel symbols vanish at $(p,t)$.

 We
compute
\begin{eqnarray}\nonumber
\frac{\partial}{\partial t}V  &  = &\frac{\partial}{\partial
t}\left( \nabla X^* \right) =-\left( \frac{\partial}{\partial
t}\Gamma_{kj}^{l}\right)  X^*_{l}dx^k\otimes dx^j +\nabla\left(
\frac{\partial}{\partial t}X^*\right) \\ \nonumber
 &= &
\left(-\nabla^{l}\Ric_{kj}+\nabla_{k}\Ric_j^l+\nabla_{j}\Ric_k^l\right)X^*_{l}dx^k\otimes
dx^j \\
& & +\nabla\left(  \Delta X^*-\Ric(X,\cdot)\right)
.\label{eqn1}
\end{eqnarray}

We have \begin{eqnarray*} \nabla(\triangle X^*) & = &
\nabla\left((g^{rs}\left(\nabla_r\nabla_s(X^*)-\Gamma_{rs}^l\nabla_lX^*\right)\right)
\\
&= &
g^{rs}\left(\nabla\left(\nabla_r\nabla_s(X^*)-\Gamma_{rs}^l\nabla_lX^*\right)\right).
\end{eqnarray*}

Let us recall the formula for commuting $\nabla$ and $\nabla_r$. The
following is immediate from the definitions.

\begin{claim}[Commuting covariant derivatives and curvature]
\label{claim:commuting-covariant-derivatives}
\uses{def:riemann-curvature-tensor}
For any tensor $\phi$ we have
$$\nabla(\nabla_r\phi)=\nabla_r(\nabla\phi)+dx^k\otimes
\Rm(\partial_k,\partial_r)(\phi)-\nabla_r(dx^l)\otimes
\nabla_l(\phi).$$
\end{claim}

Applying this to our formula gives
\begin{eqnarray*}\nabla(\triangle
X^*) & = &  g^{rs}\left(\nabla_r\nabla\nabla_sX^*+dx^k\otimes \Rm(\partial_k,\partial_r)(\nabla_sX^*)-\nabla_r(dx^l)\otimes
\nabla_l\nabla_sX^* \right. \\
& & \left. -\nabla(\Gamma_{rs}^l\nabla_lX^*)\right) .
\end{eqnarray*}
 Now we
apply the same formula to commute $\nabla$ and $\nabla_s$. The
result is
\begin{eqnarray*}
\nabla(\triangle X^*) & = & g^{rs}\left(\nabla_r\nabla_s\nabla
X^*+\nabla_r\left(dx^k\otimes \Rm(\partial_k,\partial_s)X^*-\nabla_sdx^l\otimes \nabla_lX^*\right)\right. \\
& & \left. +dx^k\otimes \Rm(\partial_k,\partial_r)(\nabla_sX^*)-\nabla_r(dx^l)\otimes
\nabla_l\nabla_sX^*-\nabla(\Gamma_{rs}^l\nabla_lX^*)\right).
\end{eqnarray*}

Now we expand
\begin{eqnarray*}
\lefteqn{\nabla_r\left(dx^k\otimes \Rm(\partial_k,\partial_s)X^*-\nabla_sdx^l\otimes
\nabla_lX^*\right)}
\\ & = & \nabla_r(dx^k)\otimes
\Rm(\partial_k,\partial_s)X^*+dx^k\otimes\nabla_r(\Rm(\partial_k,\partial_s))X^*
\\
& & + dx^k\otimes \Rm(\partial_k,\partial_s)\nabla_rX^*-\nabla_r\nabla_sdx^l\otimes\nabla_lX^*-\nabla_sdx^l\otimes
\nabla_r\nabla_lX^*
\end{eqnarray*}
Invoking the fact that the Christoffel symbols vanish at the point
of space-time where we are making the computation, this above
expression simplifies to
\begin{eqnarray*}
\lefteqn{\nabla_r\left(dx^k\otimes \Rm(\partial_k,\partial_s)X^*-\nabla_sdx^l\otimes
\nabla_lX^*\right)}
\\ & = & dx^k\otimes\nabla_r(\Rm(\partial_k,\partial_s))X^* + dx^k\otimes
\Rm(\partial_k,\partial_s)\nabla_rX^*-\nabla_r\nabla_sdx^l\otimes\nabla_lX^*.
\end{eqnarray*}

Also, expanding and using the vanishing of the Christoffel symbols
we have
\begin{eqnarray*}
-\nabla(\Gamma_{rs}^l\nabla_lX^*)& = &
-d\Gamma_{rs}^l\otimes\nabla_lX^*-\Gamma_{rs}^l\nabla\nabla_lX^* \\
& = & -d\Gamma_{rs}^l\otimes\nabla_lX^*.
\end{eqnarray*}



Plugging these computations into equation above and using once more the
vanishing of the Christoffel symbols gives
\begin{eqnarray*}
\nabla(\triangle X^*) & = & \triangle(\nabla X^*)+g^{rs}\left(
dx^k\otimes\nabla_r(\Rm(\partial_k,\partial_s))X^* +
dx^k\otimes
\Rm(\partial_k,\partial_s)\nabla_rX^* \right.\\
& & \left.
  -\nabla_r\nabla_sdx^l\otimes\nabla_lX^* +dx^k\otimes
\Rm(\partial_k,\partial_r)(\nabla_sX^*)
 -d\Gamma_{rs}^l\otimes\nabla_lX^*\right).
\end{eqnarray*}

 Now by the symmetry of $g^{rs}$ we can amalgamate the second and
 fourth terms on the right-hand side to give
\begin{eqnarray*}
\nabla(\triangle X^*) & = & \triangle(\nabla X^*)+g^{rs}\left(
dx^k\otimes\nabla_r(\Rm(\partial_k,\partial_s))X^* \right. \\
 & & \left. + 2dx^k\otimes
\Rm(\partial_k,\partial_s)\nabla_rX^*
-\nabla_r\nabla_sdx^l\otimes\nabla_lX^*
 -d\Gamma_{rs}^l\otimes\nabla_lX^*\right).
\end{eqnarray*}


We expand
$$\Rm(\partial_k,\partial_s)\nabla_rX^* =-{{R_{ks}}^{l}}_jV_{rl}dx^j.$$
Also we have (again using the vanishing of the Christoffel symbols)
\begin{eqnarray*}
-\nabla_r\nabla_sdx^l-d\Gamma_{rs}^l & = &
\nabla_r\Gamma_{ks}^ldx^k-\partial_k\Gamma_{rs}^ldx^k \\
& = & {{R_{rk}}^l}_s dx^k.
\end{eqnarray*}

Lastly,
\begin{eqnarray*}
\nabla_r(\Rm(\partial_k,\partial_s))X^*=-{{(\nabla_rR)_{ks}}^l}_jX^*_ldx^j.
\end{eqnarray*}


Plugging all this in and raising indices yields
\begin{eqnarray*}
\nabla(\triangle X^*) & = & \triangle(\nabla X^*)-
g^{rs}{{(\nabla_rR)_{ks}}^l}_jX^*_ldx^k\otimes
dx^j-2{{R_k}^{rl}}_jV_{rl}dx^k\otimes dx^j \\
& &  +
g^{rs}{{R_{rk}}^l}_sV_{lj}dx^k\otimes dx^j \\
& = & \triangle(\nabla X^*)- g^{rs}{{(\nabla_rR)_{ks}}^j}_lX^*_jdx^k\otimes
dx^l-2{{R_k}^{rl}}_jV_{rl}dx^k\otimes dx^j  \\
& & -\Ric_k^lV_{lj}dx^k\otimes dx^j.
\end{eqnarray*}

Thus, we have
\begin{eqnarray*}
\lefteqn{\nabla(\triangle X^*) -\nabla(\Ric(X,\cdot)^*)=} \\ &  &
\triangle(\nabla X^*)- g^{rs}{{(\nabla_rR)_{ks}}^l}_jX^*_ldx^k\otimes
dx^j-2{{R_k}^{rl}}_jV_{rl}dx^k\otimes dx^j
\\
& & -\left(\Ric_k^lV_{lj}+\nabla_k(\Ric)_j^lX^*_l+\Ric_j^lV_{kl}\right)dx^k\otimes dx^j,
\end{eqnarray*}
and consequently, plugging back into Equation~(\ref{eqn1}), and
canceling the two like terms appearing with opposite sign, we have
\begin{eqnarray*}
\frac{\partial}{\partial t}V  & = & \left(-\nabla^{l}\Ric_{kj}+\nabla_{j}\Ric_k^l\right)X^*_{l}dx^k\otimes
dx^j +\triangle(\nabla X^*) \\
& & -g^{rs}{{(\nabla_rR)_{ks}}^l}_jX^*_ldx^k\otimes
dx^j-2{{R_k}^{rl}}_jV_{rl}dx^k\otimes dx^j
\\
& & -\left(\Ric_k^lV_{lj}+\Ric_j^lV_{kl}\right)dx^k\otimes
dx^j.
\end{eqnarray*}


The last thing we need to see in order to complete the proof is that
$$-g^{rs}{{(\nabla_rR)_{ks}}^l}_j-\nabla^{l}\Ric_{kj}+\nabla_{j}\Ric_k^l=0.$$
This is obtained by  contracting $g^{rs}$ against the Bianchi
identity
$$\nabla_r{{R_{ks}}^l}_j+\nabla^lR_{ksjr}+\nabla_j{R_{ksr}}^l=0.$$
\end{proof}


Let $h_{ij}$ be defined by $h_{ij}= V_{ij}+V_{ji}$. It follows from
(\ref{eq.4}) that
\begin{equation}
\frac{\partial}{\partial t}h_{ij}=\Delta_{L}h_{ij}, \label{eq.5}
\end{equation}
where by definition $\Delta_{L}h_{ij}=\Delta
h_{ij}+2{{R_i}^{kl}j}h_{kl}-\Ric_i^k h_{kj}-\Ric_j^kh_{ki}$ is the \textbf{Lichnerowicz
Laplacian}. A simple calculation shows
that there is a constant $C>0$ such that
\begin{align}
\left(  \frac{\partial}{\partial t}-\Delta\right)  \left\vert h_{ij}
\right\vert ^{2}  &  =-2\left\vert \nabla_{k}h_{ij}\right\vert ^{2}
+4R^{ijkl}h_{jk}h_{il} \label{eq.6}\\
\frac{\partial}{\partial t}\left\vert h_{ij}\right\vert ^{2}  & \leq
\Delta\left\vert h_{ij}\right\vert ^{2}-2\left\vert \nabla_{k}h_{ij}
\right\vert ^{2}+C\left\vert h_{ij}\right\vert ^{2}. \label{eq.7}
\end{align}
Note that $X(t)$ is a Killing vector field for $g(t)$ if and only if
$h_{ij}(t)=0$. Since Equation~(\ref{eq1}) is linear and since the
curvature is bounded on each time-slice, for any given bounded
Killing vector field $X(0)$ for metric $g(0)$, there is a bounded
solution $X^{i}\left( t\right) $ of Equation~(\ref{eq1}) for
$t\in\lbrack0,T]$. Then $\left \vert h_{ij}\left( t\right) \right
\vert^2 $ is a bounded function satisfying (\ref{eq.7}) and $\left
\vert h_{ij} \right \vert^2 (0)=0$. One can apply the maximum
principle to (\ref{eq.7}) to conclude that $h_{ij}(t)=0$ for all
$t\geq 0$. This is done as follows: Let $h(t)$ denote the maximum of
$|h_{ij}(x,t)|^2$ on the $t$ time-slice. Note that, for any fixed
$t$ the function $|h_{ij}(x,t)|^2$ approaches $0$ as $x$ tends to
infinity  since the metric is asymptotic at infinity to the product
of a round metric on $S^2$ and the standard metric on the line. By
virtue of (\ref{eq.7}) and Proposition~\ref{prop:forward-difference-maximum}, the
function $h(t)$ satisfies $dh/dt\le Ch$ in the sense of forward
difference quotients, so that $d(e^{-Ct}h)/dt\le 0$, also in the
sense of forward difference quotients. Thus, by
Corollary~\ref{lem:forward-difference-quotient}, since $h(0)=0$ and $h\ge 0$, it follows
that $e^{-Ct}h(t)=0$ for all $t\ge 0$, and consequently, $h(t)=0$
for all $t\ge 0$.

Thus, the evolving vector field $X( t) $ is a Killing vector field
for $g(t)$ for all $t\in [0,T)$. The following is a very nice
observation of Bennett Chow; we thank him for allowing us to use it
here. From $h_{ij}=0$ we have $\nabla_{j} X^i +\nabla_{i} X^{j}=0$.
Taking the $\nabla_{j}$ derivative and summing over $j$ we get
$\Delta X^i +R^i_kX^k =0$ for all $t$. Hence (\ref{eq.3}) gives
$\frac{\partial}{\partial t}X^{i}=0$ and $X(t)=X(0)$, i.e., the
Killing vector fields are stationary and remain Killing vector
fields for the entire flow $g(t)$.
 Since at $t=0$ the Lie algebra $so(3)$ of the standard rotation action
 consists of Killing vector fields, the same is true for all the metrics $g(t)$
 in the standard solution.
 Thus, the
rotation group $SO(3)$ of $\mathbb{R}^3$ is  contained in the
isometry group of $g(t)$ for every $t\in [0,T)$. We have shown:

\begin{corollary}[Standard solution is rotationally symmetric]
\label{cor:standard-solution-rotationally-symmetric}
\uses{def:standard-ricci-flow}
The standard solution $g(t),\ t\in [0,T)$, consists of a family of metrics
all of which are
rotationally symmetric by the standard action of $SO(3)$ on
$\mathbb{R}^3$.
\end{corollary}










\subsection{Non-collapsing}



\begin{proposition}[Non-collapsing of the standard flow]
\label{prop:standard-flow-kappa-noncollapsed}
\label{stdkappa}
\uses{def:standard-ricci-flow}
For any $r>0$ sufficiently small, there is a $\kappa>0$ such that
the standard flow is $\kappa$-non-collapsed on all scales $\le r$.
\end{proposition}


\begin{proof}
\uses{prop:forward-difference-maximum,n/2,THM}
Since the curvature of the standard solution is non-negative, it
follows directly that $2|\Ric|^2\le R^2$. By
Equation~(\ref{Revol}) this gives
$$\frac{\partial R}{\partial t}=\triangle R+2|\Ric|^2\le \triangle
R+R^2.$$ Let $C=\max(2,\max_{x\in \Ar^3}R(x,0))$. Suppose
that $t_0< T$ and $t_0<1/C$.

\begin{claim}[Curvature bound along the standard flow]
\label{claim:standard-flow-curvature-bound}
\uses{def:standard-ricci-flow}
For all $x\in \Ar^3$ and $t\in [0,t_0]$ we have
$$R(x,t)\le \frac{C}{1-Ct}.$$
\end{claim}


\begin{proof}
\uses{prop:standard-flow-asymptotic-cylinder}
By the asymptotic condition, there is a compact subset $X\subset
\Ar^3$ such that for any point $p\in \Ar^3\setminus X$ and for any
$t\le t_0$ we have $R(p,t)<2/(1-t)$. Since $C\ge 2$, for all $t$ for
which $\sup_{x\in \Ar^3}R(x,t)\le 2/(1-t)$,  we also have
$$R(x,t)\le \frac{C}{1-Ct}.$$

Consider the complementary subset of $t$, that is to say the subset
of $[0,t_0]$ for which there is $x\in \Ar ^3$ with
$R(x,t)>C/(1-Ct)$. This is an open subset of $[0,t_0]$, and hence is
a disjoint union of relatively open intervals. Let $\{t_1<t_2\}$ be
the endpoints of one such interval. If $t_1\not=0$, then clearly
$R_{\max}(t_1)=C/(1-Ct_1)$. Since $C\ge \sup_{x\in
\Ar^3}R(x,0)$, this is also true if $t_1=0$. For every  $t\in
[t_1,t_2]$ the maximum of $R$ on the $t$ time-slice is achieved, and
the subset of $\Ar^3\times [t_1,t_2]$ of all points where maxima are
achieved is compact. Furthermore, at any maximum point we have
$\partial R/\partial t\le R^2$. Hence, according to
Proposition~\ref{prop:forward-difference-maximum} for all $t\in [t_1,t_2]$ we have
$$R_{\max}(t)\le G(t)$$
where $G'(t)=G^2(t)$ and $G(t_1)=C/(1-Ct_1)$. It is easy to see that
$$G(t)= \frac{C}{1-Ct}.$$
This shows that for all $t\in [t_1,t_2]$ we have $R(x,t)\le
\frac{C}{1-Ct}$, completing the proof of the claim.
\end{proof}


This shows that for $t_0<T$ and $t_0<1/C$ the scalar curvature is
bounded on $M\times [0,t_0]$ by a constant depending only on $C$ and
$t_0$. Since we are assuming that our flow is maximal, it follows
that $T\ge 1/C$.


Since  $(\Ar^3,g_0)$ is asymptotic to $(S^2\times \Ar,h(0)\times
ds^2)$, by compactness there is $V>0$ such that for any metric ball
$B(x,0,r)$ on which $|\Rm|\le r^{-2}$ we have $\Vol\,B(x,r)\ge V r^3$. Since there is a uniform bound on the
curvature on $[0,1/2C]$, it follows that there is $V'>0$ so that any
ball $B(q,t,r)$ with $t\le 1/2C$ on which $|\Rm|\le r^{-2}$
satisfies $\Vol\,B(q,t,r)\ge V'r^3$. Set $t_0=1/4C$. For any
point $x=(p,t)$ with $t\ge 1/2C$ there is a point $(q,t_0)$ such
that $l_x(q,t_0)\le 3/2$; this by Theorem~\ref{n/2}. Since
$B(q,0,1/\sqrt{R_{\max}(0)})\subset \Ar^3$ has volume at least
$V/R_{\max}(0)^{3/2}$, and clearly $l_x$ is bounded above on
$B(q,0,1/\sqrt{R_{\max}(0)})$ by a uniform constant, we see that
the reduced volume of $B(q,0,1/\sqrt{R_{\max}(0)})$ is uniformly
bounded from below. It now follows from Theorem~\ref{THM} that there
is $\kappa_0>0$ such that if $|\Rm|$ is bounded by $r^{-2}$ on
the parabolic neighborhood $P(p,t,r,-r^2)$ and $r\le \sqrt{1/4C}$,
then the volume of this neighborhood is at least $\kappa_0r^3$.
Putting all this together we see that there is a universal
$\kappa>0$ such that the standard solution is $\kappa$-non-collapsed
on all scales at most $\sqrt{1/4C}$.
\end{proof}

\section{Uniqueness}\label{12.4}

Now we turn to the proof of uniqueness. The idea is to mimic the
proof of uniqueness in the compact case, by replacing the Ricci flow
by a strictly parabolic flow. The material we present here is
closely related to and derived from the presentation given in
\cite{LuTian}. The presentation here is the analogy in the context
of the standard solution of DeTurck's argument presented in
Section~\ref{3.3}.

\subsection{From Ricci flow to Ricci-DeTurck flow}

In this subsection we discuss the Ricci-DeTurck flow and the
harmonic map flow. Let $(M^n,g(t)),t \in [t_0,T]$ be a solution of
the Ricci flow and let $\psi_{t}\colon M \rightarrow M,\ t \in
\left[ t_0,T_{1}\right]$ be a solution of the harmonic map
flow
\begin{align}
\frac{\partial\psi_{t}}{\partial t}  &  =\Delta_{g\left(  t\right)
,g(t_0) }\psi_{t},\label{eq2.1} \quad \psi_{t_0}  =\Id.
\end{align}

Here, $\Delta_{g\left(  t\right) ,g(t_0) }$ is the Laplacian for
maps from the Riemannian manifold $(M,g(t))$ to the Riemannian
manifold $(M,g(t_0))$. In local coordinates $(x^i)$ on the domain
$M$ and $(y^\alpha)$ on the target $M$, the harmonic map flow
(\ref{eq2.1}) can be written as
\begin{equation}
\left( \frac{\partial}{\partial t}-\Delta_{g\left(  t\right)
}\right) \psi^{\alpha}\left(  x,t\right)  =g^{ij}\left( x,t\right)
\Gamma _{\beta\gamma}^{\alpha}\left(  \psi\left(  x,t\right) \right)
\frac {\partial\psi^{\beta}\left(  x,t\right)  }{\partial
x^{i}}\frac{\partial \psi^{\gamma}\left(  x,t\right)  }{\partial
x^{j}} \label{eq2.2}
\end{equation}
where $\Gamma_{\beta\gamma}^{\alpha}$ are the Christoffel symbols of
$g(t_0)$. Suppose $\psi\left(x, t\right)$ is a bounded smooth
solution of~\ref{eq2.2} with $\psi_{t_0}=\Id$. Then $\psi\left(
t\right),\ t \in[t_0, T_1]$ are diffeomorphisms when $T_{1}>t_0$ is
sufficiently close to $t_0$. For any such $T_1$ and for $t_0 \leq t
\leq T_1$, define $\hat{g}\left( t\right) =\left(
\psi_{t}^{-1}\right) ^{\ast}g\left( t\right)$. Then $\hat g(t)$
satisfies the following equation:
\begin{equation}\label{RDfloweqn}
\frac{\partial}{\partial t}\hat{g}_{ij} =-2\widehat{\Ric}_{ij}+\hat{\nabla} _{i}W_{j}(t)+\hat{\nabla}_{j}W_{i}(t) \quad
\quad \hat{h}\left( 0\right) =h(0),
\end{equation}
where $\widehat{\Ric}_{ij}$ and $\hat{\nabla}_{i}$ are the
Ricci curvature and Levi-Civita connection of $\hat{g}(t)$
respectively and $W(t)$ is the time-dependent $1$-form defined by
\[
W(t) _{j}=\hat{g}_{jk}(t)\hat{g}\,^{pq}(t)\left(
\hat{\Gamma}_{pq}^{k}(t)-\Gamma_{pq}^{k}(t_0)   \right).
\]
Here, $\hat{\Gamma}_{pq}^{k}(t)$ denotes the Christoffel symbols of
the metric $\hat g(t)$ and $\Gamma_{pq}^k(t_0)$ denotes the
Christoffel symbols of the metric $g(t_0)$. (See, for example,
(\cite{Shi1} Lemma 2.1).)
 We call a solution to this flow equation a
 \textbf{Ricci-DeTurck flow} (see \cite{DeTurck}, or \cite{ChowKnopf}
Chapter 3 for details). In local coordinates we have

\begin{align}
& \frac{\partial\hat{g}_{ij}}{\partial t}  =\hat{g}^{kl}\nabla
_{k}\nabla_{l}\hat{g}_{ij}-\hat{g}^{kl}g(t_0)_{ ip}\hat{g}^{pq}R_{jk
ql}\left(  g(t_0)\right)  -\hat{g} ^{kl}g(t_0)_{
jp}\hat{g}^{pq}R_{ik ql}\left(  g(t_0)\right) \nonumber
\\
&  +\frac{1}{2}\hat{g}^{kl}\hat{g}^{pq}\left[
\begin{array}
[c]{c} \nabla_{i}\hat{g}_{pk}\nabla_{ j}\hat{g}_{ql}+2\nabla
_{k}\hat{g}_{jp}\nabla_{ q}\hat{g}_{il}\\
-2\nabla_{k}\hat{g}_{jp}\nabla_{l}\hat{g}_{iq}-2\nabla_{
j}\hat{g}_{pk}\nabla_{l}\hat{g}_{iq}-2\nabla_{ i}\hat
{g}_{pk}\nabla_{l}\hat{g}_{jq}
\end{array}
\right]  . \label{eqRcDeTurck}
\end{align}
where $\nabla$ is the Levi-Civita connection of $g(t_0)$. This is a
strictly parabolic equation.


\begin{lemma}[Harmonic map flow gives a Ricci-DeTurck flow]
\label{lem:harmonic-map-flow-ricci-deturck}
\label{12.7}
\uses{def:standard-ricci-flow}
Suppose that $g(t)$ solves the Ricci flow equation and suppose that
$\psi_t$ solves the harmonic map flow equation,
Equation~(\ref{eq2.1}); then $\hat g(t)=(\psi_t^{-1})^*g(t)$ solves
the Ricci-DeTurck flow, Equation~(\ref{RDfloweqn}) and $\psi_t$
satisfies the following ODE:
$$\frac{\partial \psi_t}{\partial t}=-\hat g^{ij}(t)W(t).$$
\end{lemma}

\begin{proof}
The first statement follows from the second statement and a standard
Lie derivative computation. For the second statement, we need to
show
$$\triangle_{g(t),g(0)} \psi^\alpha = -\hat
g^{pq}\left(\hat\Gamma_{pq}^\alpha(t)-\Gamma_{pq}^\alpha(t_0)\right).$$
Notice that this equation is a tensor equation, so that we can
choose coordinates in the domain and range so that $\Gamma(t)$
vanishes at the point $p$ in question and $\Gamma(t_0)$ vanishes at
$\psi_t(p)$.  With these assumptions we need to show
$$g^{pq}(t)\frac{\partial^2 \psi^\alpha}{\partial x^p\partial x^q}=-\hat g^{pq}(t)\hat\Gamma_{pq}^\alpha(t).$$
This is a direct computation using the change of variables formula
relating $\hat \Gamma$ and $\Gamma$.
\end{proof}

\begin{corollary}[Ricci-DeTurck uniqueness implies Ricci flow uniqueness]
\label{cor:ricci-deturck-uniqueness-implies-ricci-uniqueness}
\label{dtuniqtouniq}
\uses{lem:harmonic-map-flow-ricci-deturck}
Suppose that $(M,g_1(t)),\ t_0\le t\le T$, and $(M,g_2(t)),\ t_0\le
t\le T$, are solutions to the Ricci flow equation for which there
are solutions
$$\psi_{1,t}\colon (M,g_1(t))\to (M,g_1(0))$$ and
$$\psi_{2,t}\colon
(M,g_2(t))\to (M,g_2(0))$$ to the harmonic map equation with
$\psi_{1,t_0}=\psi_{2,t_0}=\Id$. Let $\hat
g_1(t)=(\psi_{1,t}^{-1})^*g_1(t)$ and $\hat
g_2(t)=(\psi_{2,t}^{-1})^*g_2(t)$ be the corresponding solutions to
the Ricci-DeTurck flow. Suppose that $\hat g_1(t)=\hat g_2(t)$ for
all $t\in [t_0,T]$. Then $g_1(t)=g_2(t)$ for all $t\in [t_0,T]$.
\end{corollary}

\begin{proof}
Since $\psi_{a,t}$ satisfies the equation
$$\frac{\partial \psi_{a,t}}{\partial t}=-\hat g_a^{ij}W(t)_j$$
where the time-dependent vector field $W(t)$ depends only on $\hat
g_a$, we see that $\psi_{1,t}$ and $\psi_{2,t}$ both solve the same
time-dependent ODE and since $\psi_{1,t_0}=\psi_{2,t_0}=\Id$,
it follows that $\psi_{1,t}=\psi_{2,t}$ for all $t\in [t_0,T]$. On
the other hand, $g_a(t)=\psi_{a,t}^*\hat g_a(t)$, so that it follows
that $g_1(t)=g_2(t)$ for $t\in [t_0,T]$.
\end{proof}


Our strategy of proof is to begin with a standard solution $g(t)$
and show that there is a solution to the harmonic map equation for
this Ricci flow with appropriate decay conditions at infinity. It
follows that the solution to the Ricci-DeTurck flow constructed is
well-controlled at infinity. Suppose that we have two standard
solutions $g_1(t)$ and $g_2(t)$ (with the same initial conditions
$g_0$) that agree on the interval $[0,t_0]$ which is a proper
subinterval of the intersection of the intervals of definition of
$g_1(t)$ and $g_2(t)$. We construct solutions to the harmonic map
flow equation from $g_a(t)$ to $g_a(t_0)$ for $a=1,2$. We show that
solutions always exist for some amount of time past $t_0$. The
corresonding Ricci-DeTurck flows $\hat g_a(t)$ starting at $g_{t_0}$
are well-controlled at infinity. Since the Ricci-DeTurck flow
equation is a purely parabolic equation, it has a unique solution
with appropriate control at infinity and given initial condition
$g_1(t_0)=g_2(t_0)$. This implies that the two Ricci-DeTurck flows
we have constructed are in fact equal. Invoking the above corollary,
we conclude that $g_1(t)$ and $g_2(t)$ agree on a longer interval
extending  past $t_0$. From this it follows easily that $g_1(t)$ and
$g_2(t)$ agree on their common domain of definition. Hence, if they
are both maximal flows, they must be equal.



\section{Solution of the harmonic map flow}

In order to pass from a solution to the Ricci flow equation to a
solution of the Ricci-DeTurck flow we must prove the existence of a
solution of the harmonic map flow
associated with the Ricci flow. In this section we study the
existence of the harmonic flow (\ref{eq2.1}) and its asymptotic
behavior at the space infinity when $h(t)=g(t)$ is a standard
solution. Here we use in an essential way the rotationally symmetric
property and asymptotic property at infinity of $g(t)$. In this
argument there is no reason, and no advantage, to restricting to
dimension three, so we shall consider rotationally symmetric
complete metrics on $\Ar^n$, i.e., complete metrics on $\Ar^n$
invariant under the standard action of $SO(n)$. Let
$\theta=(\theta^1,\cdots,\theta^{n-1})$ be local coordinates on the
round $(n-1)$-sphere of radius $1$, and let $d\sigma$ be the metric
on the sphere. We denote by $\hat r$ the standard radial coordinate
in $\Ar^n$. Since $g(t)$ is rotationally symmetric and $n \geq 3$,
we can write

\begin{align}
g(t)=dr^2+f(r,t)^2 d \sigma \label{eq2.4}
\end{align}
Here $r=r(\hat r,t)$ is the (time-dependent) radial coordinate on
$\mathbb{R}^n$ for the metric $g(t)$.

\begin{claim}[The radial coordinate $r$ as a function of $\hat r$ and $t$]
\label{claim:radial-coordinate-properties}
\label{cl12.20}
\uses{def:standard-ricci-flow}
For any fixed $t$ the function $r\colon \Ar^n \to[0,\infty)$ is a
function only of $\hat r$. Considered as a function of two
variables,
 $r(\hat r,t)$ is a smooth function defined for $\hat r\ge 0$. It is
 an odd function of $\hat r$. For fixed $t$ it is an increasing
 function of $\hat r$.
\end{claim}



\begin{proof}
Write the metric $g(t)=g_{ij}dx^idx^j$ and let
\begin{equation}
x^1  =  \hat{r} \cos \theta^1, \ \ x^2 =  \hat{r} \sin \theta^1 \cos
\theta^2, \cdots,
 x^n  =  \sin \theta^1 \cdots \sin \theta^{n-1}.
 \end{equation}
 We compute $f(r,t)$
by restricting attention to the ray $\hat{r}=x^1$ and $\theta^1
=\cdots=\theta^{n-1}=0$, i.e., $x^2 =\cdots =x^n=0$. Then
\[
g(t)=g_{11}(\hat{r},0,\cdots,0,t)d\hat{r}^2+g_{22}(\hat{r},0,\cdots,0,t)
\hat{r}^2 d \sigma.
\]
Both $g_{11}$ and $g_{22}$ are positive smooth and even in $\hat r$.
Clearly,$\sqrt{g_{11}(\hat r,0,\cdots,0,t)}$ is a positive smooth
function defined for all $(\hat r,t)$ and is invariant under the
involution $\hat r\mapsto -\hat r$. Hence its restriction to $\hat
r\ge0$ is an even function. Since
\[
r= \int_0^{\hat{r}} \sqrt{g_{11}(\hat{s},0,\cdots,0,t)}d\hat{s} =
\hat{r} \int_0^{1} \sqrt{g_{11}(\hat{r}s,0,\cdots,0,t)}ds,
\]
we see that $r(\hat r,t)$ is of the form $\hat r\cdot \phi(\hat
r,t)$ where $\phi(\hat r,t)$ is an even smooth function. This shows
that $r(\hat r,t)$ is an odd function. It is also clear from this
formula that $\partial r/\partial \hat r>0$.
\end{proof}




 Since, for each $t_0$, the function $r(\hat
r,t_0)$ is an increasing function of $\hat r$,  it can be inverted
to give a function $\hat r(r,t_0)$. In Equation~(\ref{eq2.4}), we
have chosen to write $f$ as a function of $r$ and $t$, rather than a
function of $\hat r$ and $t$. We look for rotationally symmetric
solutions to Equation~(\ref{eq2.1}), i.e., solutions of the form:

\begin{align}\label{ansatz}
\psi(t)\colon \mathbb{R}^n \rightarrow  \mathbb{R}^n \qquad
\psi(t)(r,\theta) & =(\rho(r,t),\theta) \ \ \ \text{for}\ \ t\ge t_0 \\
\psi(r,t_0) & = \Id \nonumber
\end{align}

We shall adopt the following conventions: we shall consider
functions $f(w,t)$ defined in the closed half-plane $w\ge 0$. When
we say that such a function is {\em smooth} we mean that for each
$n,m\ge 0$ we have a continuous function $f_{nm}(w,t)$ defined for
all $w\ge 0$ with the following properties:
\begin{enumerate}
\item[(1)] $$f_{00}=f$$
\item[(2)] $$\frac{\partial f_{nm}}{\partial t}=f_{n(m+1)}$$
\item[(3)] $$\frac{\partial f_{nm}}{\partial w}=f_{(n+1)m},$$
\end{enumerate}
where in Item (3) the partial derivative along the boundary $w=0$ is
a right-handed derivative only. We say such a function is {\em even}
if $f_{(2k+1)m}(0,t)=0$ for all $k\ge 0$.

We have the following elementary lemma:

\begin{lemma}[Even and odd radial functions extend smoothly]
\label{lem:even-odd-radial-functions}
(a) Suppose that $f(w,t)$ is a smooth function defined for $w\ge 0$.
Define $\phi(r,t)=f(r^2,t)$. Then $\phi(r,t)$ is a smooth function
defined for all $r\in \Ar$. Now fix $k$ and let $\hat r\colon
\Ar^k\to [0,\infty)$ be the usual radial coordinate.  Then we have a
smooth family of smooth functions on $\Ar^k$ defined by
$$\hat\phi(x^1,\ldots,x^k,t)=\phi(\hat r(x^1,\ldots,x^k),t)=f(\sum_{i=1}^k(x^i)^2,t).$$

\noindent (b) If $\psi(r,t)$ is a smooth function defined for $r\ge
0$ and if it is even in the sense that its Taylor expansion to all
orders along the line $r=0$ involves only even powers of $r$, then
there is a smooth function $f(w,t)$ defined for $w\ge 0$ such that
$\psi(r,t)=f(r^2,t)$. In particular, for any $k\ge 2$ the function
$\hat\psi((x^1,\ldots,x^k),t)=\psi(r(x^1,\ldots,x^k),t)$ is a smooth
family of smooth functions on $\Ar^k$.
\end{lemma}

\begin{proof}
Item (a) is obvious, and Item (b) is obvious away from $r=0$. We
establish Item (b) along the line $r=0$. Consider the Taylor theorem
with remainder to order $2N$ in the $r$-direction for $\psi(r,t)$ at
a point $(0,t)$. By hypothesis it takes the form
$$\sum_{i=0}c_i(t)w^{2i}+w^{2N+1}R(w,t).$$
Now we replace $w$ by $\sqrt{r}$ to obtain
$$f(r,t)=\sum_{i=0}c_ir^i+\sqrt{r}^{2N+1}R(\sqrt{r},t).$$
Applying the usual chain rule and taking limits as $r\rightarrow
0^+$ we see that $f(r,t)$ is $N$ times differentiable along the line
$r=0$. Since this is true for every $N<\infty$, the result follows.
\end{proof}


Notice that an even function $f(r,t)$ defined for $r\ge 0$ extends
to a smooth function on the entire plane invariant under $r\mapsto
-r$. When we say a function $f(r,t)$ {\em defines a smooth family of
smooth functions on $\Ar^n$} we mean that, under the substitution
$\hat f((x^1,\ldots,x^n),t)=f(r(x^1,\ldots,x^n),t)$, the function
$\hat f$ is a smooth function on $\Ar^n$ for each $t$.

We shall also consider odd functions $f(r,t)$, i.e., smooth
functions defined for $r\ge 0$ whose Taylor expansion in the
$r$-direction along the line $r=0$ involves only odd powers of $r$.
These do not define smooth functions on $\Ar^n$. On the other hand,
by the same argument as above with the Taylor expansion one sees
that they can be written as $rg(r,t)$ where $g$ is even, and hence
define  smoothly varying families of smooth functions on $\Ar^n$.
Notice also that the product of two odd functions $f_1(r,t)f_2(r,t)$
is an even function and hence this product defines a smoothly
varying family of smooth function on $\Ar^n$.


\subsection{The properties of $r$ as a function of $\hat r$ and $t$}

We shall make a change of variables and write the harmonic map flow
equation in terms of $r$ and $\theta$. For this we need some basic
properties of $r$ as a function of $\hat r$ and $t$. Recall that we
are working on $\Ar^n$ with its usual Euclidean coordinates
$(x^1,\ldots,x^n)$. We shall also employ spherical coordinates $\hat
r,\theta^1,\ldots,\theta^{n-1}$. (We denote the fixed radial
coordinate on $\Ar^n$ by $\hat r$ to distinguish it from the varying
radial function $r=r(t)$ that measures the distance from the tip in
the metric $g(t)$.)

As a corollary of Claim~\ref{claim:radial-coordinate-properties} we have:

\begin{corollary}[$r^2$ is a smooth function of the ambient coordinates]
\label{cor:r-squared-smooth-function}
\label{rcor}
\uses{claim:radial-coordinate-properties,lem:even-odd-radial-functions}
$r^2(\hat r,t)$ is a smoothly varying family of smooth functions on
$\Ar^n$. Also, $\hat r$ is a smooth  function of $(r,t)$ defined for
$r\ge 0$ and odd in $r$. In particular, any smooth even function of
$r$ is a smooth even function of $\hat r$ and thus defines a smooth
function on $\Ar^n$. Moreover, there is a smooth function $\xi(w,t)$
such that $d(\log\, r)/dt=r^{-1}(dr/dt)=\xi(r^2,t)$.
\end{corollary}

For future reference we define
\begin{equation}\label{Bdefn}
B(w,t)=\frac{1}{2}\int_0^w\xi(u,t)du. \end{equation} Then $B(r^2,t)$
is a smooth function even in $r$ and hence, as $t$ varies, defines a
smoothly varying family of smooth functions on $\Ar^n$. Notice that
$$\frac{\partial B(r^2,t)}{\partial r}=2r\frac{\partial B}{\partial
w}(w,t)\left|_{w=r^2}\right.=2r\left(\frac{1}{2}\xi(r^2,t)\right)=\frac{dr}{dt}.$$

Now let us consider $f(r,t)$.

\begin{claim}[$f(r,t)$ is a smooth odd function]
\label{claim:f-odd-smooth-function}
\uses{cor:r-squared-smooth-function}
$f(r,t)$ is a smooth function defined for $r\ge 0$. It is an odd
function of $r$.
\end{claim}


\begin{proof}
We have
$$f(r,t)=\hat r(r,t)\sqrt{g_{22}(\hat r(r,t),0,\ldots,0,t)}.$$
Since $\sqrt{g_{22}(\hat r,0,\ldots,0,t)}$ is a smooth function of
$(\hat r,t)$ defined for $\hat r\ge 0$ and since it is an  even
function of $\hat r$, it follows immediately from the fact that
$\hat r$ is a smooth odd function of $r$, that $f(r,t)$ is a smooth
odd function of $r$.
\end{proof}


\begin{corollary}[Factorization $f(r,t)=r\,h(r^2,t)$]
\label{cor:f-factorization-rh}
\uses{claim:f-odd-smooth-function}
There is a smooth function $h(w,t)$ defined for $w\ge 0$ so that
$f(r,t)=rh(r^2,t)$. In particular, $h(r^2,t)$ defines a smooth
function on all of $\Ar^n$. Clearly, $h(w,t)>0$ for all $w\ge 0$ and
all $t$.
\end{corollary}

We set $\tilde h(w,t)=\log(h(w,t))$, so that $f(r,t)=re^{\tilde
h(r^2,t)}$. Notice that $\tilde h(r^2,t)$ defines a smooth function
of $\hat r^2$ and $t$ and hence is a smoothly varying family of
smooth functions on $\Ar^n$.



\subsection{The harmonic map flow equation}.

Let $\psi(t)\colon \Ar^n\to \Ar^n$ be a smoothly varying family of
smooth functions as given in Equation~(\ref{ansatz}). Using
(\ref{eq2.4}) and (\ref{ansatz}) it is easy to calculate the energy
functional using spherical coordinates with $r$ as the radial
coordinate.



\begin{align*}
E(\psi(t)) & = \frac{1}{2} \int_{\mathbb{R}^n} \left \vert \nabla
\psi(t)
\right \vert^2_{g(t),g_{t_0}} dV_{g(t)}  \\
&= \frac{1}{2} \int_{\mathbb{R}^n} \left[ \left( \frac{\partial
\rho}{\partial r} \right)^2 + (n-1) f^2(\rho,t_0) f^{-2}(r,t)
\right] dV_{g(t)}.
\end{align*}

If we have a compactly supported variation $\delta \rho =w$, then
letting $d\vol_\sigma$ denote the standard volume element on
$S^{n-1}$, we have
\begin{align*}
& \delta E(\psi(t))(w) = \frac{1}{2} \int_{\mathbb{R}^n} \left[2
\frac{\partial \rho}{\partial r} \frac{\partial w}{\partial r}
+2(n-1) f(\rho,t_0)
\frac{\partial f(\rho,t_0)}{\partial \rho}f^{-2}(r,t)w \right] dV_g(t) \\
&= \int_0^{+\infty} \left[  f^{n-1}(r,t)\frac{\partial
\rho}{\partial r} \frac{\partial w}{\partial r} +(n-1) f(\rho,t_0)
\frac{\partial f(\rho,t_0)}{\partial
\rho}f^{n-3}(r,t)w \right] dr \cdot \int_{S^{n-1}} d\vol_\sigma \\
& = \int_0^{+\infty} \left[-\frac{\partial}{\partial
r}\left(f^{n-1}\frac{\partial\rho}{\partial r}\right)w+(n-1)
f(\rho,t_0) \frac{\partial f(\rho,t_0)}{\partial \rho}f^{n-3}(r,t)w
\right] dr \cdot \int_{S^{n-1}} d\vol_\sigma\\ &=
\int_{\mathbb{R}^n}\left[ - f^{1-n} \frac{\partial }{\partial r}
\left( \frac{\partial \rho}{\partial r} f^{n-1}\right) +(n-1)
f(\rho,t_0) \frac{\partial f(\rho,t_0)}{\partial \rho}f^{-2}(r,t)
\right] w dV_{g(t)}.
\end{align*}



The usual argument shows that for a compactly supported variation
$w$ we have
$$\delta_w\left(\frac{1}{2}\int_{\Ar^n}|\nabla_{g(t),g(t_0)}\psi|^2d\vol=\int_{\Ar^n}\langle
w,-\triangle_{g(t),g(t_0)}\psi\rangle d\vol\right).$$ Thus,
$$\triangle_{g(t),g(t_0)}\psi= \left[
f^{1-n}\frac{\partial }{\partial r} \left( \frac{\partial
\rho}{\partial r} f^{n-1}\right) -(n-1) f(\rho,t_0) \frac{\partial
f(\rho,t_0)}{\partial \rho}f^{-2}(r,t) \right]
\frac{\partial}{\partial r}$$ where we have written this expression
using the coordinates $(r,\theta)$ on the range $\Ar ^n$ (rather
than the fixed coordinates $(\hat r,\theta)$).

Now let us compute $\partial \psi/\partial t(\hat r,t)$ in these
same coordinates. (We use $\hat r$ for the coordnates for $\psi$ in
the domain to emphasize that this must be the time derivative at a
fixed point in the underlying space.) Of course, by the chain  rule,
\begin{eqnarray*}
\frac{\partial \psi(\hat r,t)}{\partial t} & = & \frac{\partial
\psi(r,t)}{\partial r}\frac{\partial r}{\partial t}+\frac{\partial
\psi(r,t)}{\partial t} \\
& = & \frac{\partial \rho(r,t)}{\partial r}\frac{\partial r(\hat
r,t)}{\partial t}+\frac{\partial \rho(r,t)}{\partial
t}.\end{eqnarray*}



Consequently,   for rotationally symmetric maps as in
Equation~(\ref{ansatz}) the harmonic map flow equation (\ref{eq2.1})
has the following form:

\begin{equation*}
 \frac{\partial \rho}{\partial t}+\frac{\partial \rho}{\partial r}\frac{\partial r}{\partial t}
 = \frac{1}{f^{n-1}(r,t)}
\frac{\partial }{\partial r} \left(  f^{n-1}(r,t) \frac{\partial
\rho}{\partial r} \right)-(n-1)f^{-2}(r,t) f(\rho,t_0)
\frac{\partial f(\rho,t_0)}{\partial \rho}
\end{equation*}
or equivalently

\begin{equation}\label{eqrothar}
\frac{\partial \rho}{\partial t}=\frac{1}{f^{n-1}(r,t)}
\frac{\partial }{\partial r} \left(  f^{n-1}(r,t) \frac{\partial
\rho}{\partial r} \right)-(n-1)f^{-2}(r,t) f(\rho,t_0)
\frac{\partial f(\rho,t_0)}{\partial \rho} -\frac{\partial
\rho}{\partial r}\frac{\partial r}{\partial t}
\end{equation}


The point of  rewriting the harmonic map equation in this way is to
find an equation for the functions $\rho(r,t), f(r,t)$ defined on
$r\ge 0$. Even though the terms in this rewritten equation involve
odd functions of $r$, as we shall see, solutions to these equations
will be even in $r$ and hence will produce a smooth solution to the
harmonic map flow equation on $\Ar ^n$.

\subsection{An equation equivalent to the harmonic map flow
equation}

We will solve (\ref{eqrothar})  for solutions of the form

\[ \rho(r,t)= re^{\tilde{\rho}(r,t)}, \ \ \ t\ge t_0;\ \ \ \tilde \rho(r,t_0)=0.
\]

For $\psi$ as in Equation~(\ref{ansatz}) to define a diffeomorphism,
it must be the case that $\rho(r,t)$ is a smooth function for $r\ge
0$ which is odd in $r$. It follows from the above expression that
$\tilde\rho(r,t)$ is a smooth function of $r$ and $t$ defined for
$r\ge 0$ and even in $r$, so that it defines a smoothly varying
family of smooth functions on $\Ar^n$. Then some straightforward
calculation shows that (\ref{eqrothar}) becomes

\begin{eqnarray}\label{newform}
\frac{\partial\tilde{\rho} }{\partial t} & = &
\frac{\partial^{2}\tilde{\rho}}{\partial r^{2}}+\frac{n+1}{r}\frac
{\partial\tilde{\rho}}{\partial r}
 +\left(  n-1\right) \frac{\partial\tilde{h}(
r^{2},t)}{\partial r}  \frac{\partial\tilde{\rho} }{\partial r}
+\left( \frac{\partial\tilde{\rho}
}{\partial r}\right)  ^{2}    \\
& & +\frac{n-1}{r^{2}}\left[  1-e^{2\tilde{h}\left(
\rho^{2},t_0\right) -2\tilde{h}\left(  r^{2},t\right)  }\right]
 +2\left(  n-1\right)  \frac{\partial\tilde{h}}{\partial w}\left(  r^{2},t\right) \nonumber \\
&  & -2\left(  n-1\right)  e^{2\tilde{h}\left(  \rho ^{2},t_0\right)
+2\tilde{\rho}  -2\tilde{h}\left( r^{2},t\right)
}\frac{\partial\tilde{h}}{\partial w}(\rho^2,t_0)
-\frac{2}{r}\frac{\partial r}{\partial t}-\frac{\partial r}{\partial
t} \frac{\partial\tilde{\rho} }{\partial r} .\nonumber
\end{eqnarray}

Note that from the definition, $\tilde{h}\left( 0,t \right)=0$, we
can write $\tilde{h}\left( w,t \right)=w \tilde{h}^*(w,t)$ where
$\tilde{h}^*(w,t)$ is a smooth function of $w\ge 0$ and $t$. So
$$
\frac{n-1}{r^{2}}\left[  1-e^{2\tilde{h}\left( \rho^{2},t_0\right)
-2\tilde{h}\left(  r^{2},t\right)  }\right] =\frac{n-1}{r^{2}}\left[
1-e^{2r^2\left[ e^{2\tilde{\rho}} \tilde{h}^*\left(
\rho^{2},t_0\right)  -\tilde{h}^*\left( r^{2},t\right) \right]
}\right]
$$
which is a smooth function of $\tilde{\rho},r^2,t$.

Let
\begin{align}
G(\tilde{\rho},w,t) = & \frac{n-1}{w}\left[ 1-e^{2\tilde{h}\left(
\rho^{2},t_0\right)  -2\tilde{h}\left( w,t\right)  }\right]
 +2\left(  n-1\right)  \frac{\partial\tilde{h}}{\partial w}\left(  w,
 t\right) \nonumber \\
& -2\left(  n-1\right)  e^{2\tilde{h}\left(  \rho ^{2},t_0\right)
+2\tilde{\rho}  -2\tilde{h}\left( w,t\right) }\frac{\partial\tilde
h}{\partial w}(\rho^2,t_0) -2\xi(w,t), \label{defofG}
\end{align}
where $\xi$ is the function from Corollary~\ref{cor:r-squared-smooth-function}. Then
$G(\tilde{\rho},w,t)$ is a smooth function defined for $w\ge 0$.
Notice that when $r$ and $\tilde \rho$ are the functions associated
with the varying family of metrics $g(t)$ and the solutions to the
harmonic map flow, then $G(\tilde\rho,r^2,t)$ defines a smoothly
varying family of smooth functions on $\Ar^n$.


We have the following form of equation (\ref{newform}):

\begin{align*}
\frac{\partial\tilde{\rho} }{\partial t} = &
\frac{\partial^{2}\tilde{\rho}}{\partial r^{2}}+\frac{n+1}{r}\frac
{\partial\tilde{\rho}}{\partial r} +\left[ ( n-1)
\frac{\partial\tilde{h}}{\partial r} -\frac{\partial B}{\partial r}
\right] \left( r^{2},t\right) \frac{\partial\tilde{\rho} }{\partial
r} +\left( \frac{\partial\tilde{\rho} }{\partial r}\right)  ^{2}
+G(\tilde\rho,r^2,t).
\end{align*}



Now we think of $\tilde{\rho}$ as a rotationally symmetric function
defined on $\mathbb{R}^{n+2}$ and let $\widehat
G(\tilde\rho,(x^1,\ldots,x^{n+2}),t)=
G(\tilde\rho,\sum_{i=1}^{n+2}(x^i)^2,t)$ and then the above equation
can be written as

\begin{align}
\frac{\partial\tilde{\rho} }{\partial t}
 = \Delta \tilde{\rho} +  \nabla [ (n-1)\tilde{h} -B] \cdot \nabla\tilde{\rho}
+ \left | \nabla \tilde{\rho} \right |^2 +G(\tilde{\rho},x,t)
\label{betterharMapeq}
\end{align}
where $\nabla$ and $\Delta$ are the Levi-Civita connection and
Laplacian defined by the Euclidean metric on  $\mathbb{R}^{n+2}$
respectively and where $B$ is the function defined in
Equation~(\ref{Bdefn}).

\begin{remark}[Why we work on $\mathbb{R}^{n+2}$]
\label{rem:choice-of-dimension-n-plus-2}
The whole purpose of this rewriting of the PDE for $\tilde \rho$ is
to present this equation in such a form that all its coefficients
represent smooth functions of $\hat r$ and $t$ that are even in
$\hat r$ and hence define smooth functions on Euclidean space of any
dimension. We have chosen to work on $\Ar^{n+2}$ because the
expression for the Laplacian in this dimension has the term
$((n+1)/r)\partial \tilde \rho/\partial r$.
\end{remark}

It is important to understand the asymptotic behavior of our
functions at spatial infinity.


\begin{claim}[Asymptotics of the harmonic map flow coefficients]
\label{claim:harmonic-map-flow-asymptotics}
\label{asympt}
\uses{prop:standard-flow-asymptotic-cylinder}
For any fixed $t$ we have the following asymptotic expansions at
spatial infinity.
\begin{enumerate}
\item[(1)]
 $e^{\tilde{h}(r^2,t)}$ is asymptotic to $\frac{1}{(1-t)r}$.
\item[(2)] $\tilde{h}(r^2,t)$ is asymptotic to $-\log r$.
\item[(3)] $\frac{\partial \tilde{h}}{\partial w}(r^2,t)$ is asymptotic to
 $-\frac{1}{2r^2}$.
\item[(4)]  $r^{-1} \frac{\partial r}{\partial t}$ is asymptotic to
$\frac{C}{r}$.
\item[(5)] $\frac{\partial B(r^2,t)}{\partial r}$ is asymptotic to  $C$.
\item[(6)] $|G( \tilde{\rho},r^2,t)| \leq C_*< \infty$ where
$C_* =C_*\left(\sup\{|\tilde{\rho}|,\tilde{h}\}\right)$ is a
constant depending only on $\sup\{|\tilde{\rho}|,\tilde{h}\}$.
\end{enumerate}
\end{claim}


\begin{proof}
\uses{prop:standard-flow-asymptotic-cylinder}
The first item is immediate from  Proposition~\ref{prop:standard-flow-asymptotic-cylinder}. The
second and third follow immediately from the first. The fourth is a
consequence of the fact that by Proposition~\ref{prop:standard-flow-asymptotic-cylinder} $dr/dt$ is
asymptotic to a constant at infinity on each time-slice. The fifth
follows immediately from the fourth and the definition of
$B(r^2,t)$. Given all these asymptotic expressions, the last is
clear from the expression for $G$ in terms of $\tilde \rho$, $r^2$,
and $t$.
\end{proof}



\subsection{The short time existence}

The purpose of this subsection is to prove the following short-time
existence theorem for the harmonic map flow equation.

\begin{proposition}[Short-time existence for the harmonic map flow]
\label{prop:harmonic-map-flow-short-time-existence}
\uses{def:standard-ricci-flow}
For any $t_0\ge 0$ for which there is a standard solution $g(t)$
defined on $[0,T_1]$ with $t_0<T_1$ there is $T>t_0$ and a solution
to Equation~(\ref{betterharMapeq}) with initial condition
$\tilde{\rho}(r,t_0)=0$  defined on the time-interval $[t_0,T]$.
\end{proposition}


 At this point to simplify the notation we shift time by $-t_0$ so
that our initial time is $0$, though our initial metric is not $g_0$
but rather is the time $t_0$-slice of the standard solution we are
considering, so that now $t_0=0$ and our initial condition is
$\tilde \rho(r,0)=0$.

Let $x=(x^1,\cdots,x^{n+2})$ and $y=(y^1,\cdots,y^{n+2})$ be two
points in $\mathbb{R}^{n+1}$ and

\[
H(x,y,t)=\frac{1}{(4 \pi t)^{(n+2)/2}} e^{-\frac{|x-y|^2}{4t}}
\]
be the heat kernel. We solve (\ref{betterharMapeq}) by successive approximation
\cite{LiTam}.

Define
$$F(x, \tilde{\rho},\nabla\tilde{\rho},t)
=  \nabla \left[ (n-1)\tilde{h}-B \right] \cdot \nabla\tilde{\rho} +
\left | \nabla \tilde{\rho} \right |^2 +G( \tilde{\rho},x,t)
$$

Let $\tilde{\rho}_0(x,t)=0$ and for $i \geq 1$ we define
$\tilde{\rho}_i$ by

\begin{align}
\tilde{\rho}_i=\int_0^t \int_{\mathbb{R}^{n+2}}
H(x,y,t-s)F(y,\tilde{\rho}_{i-1}, \nabla\tilde{\rho}_{i-1} ,t)dyds
\label{rhoIntg}
\end{align}
which solves

\begin{align}
\frac{\partial\tilde{\rho}_i}{\partial t} = \Delta \tilde{\rho}_i +
F(x, \tilde{\rho}_{i-1},\nabla\tilde{\rho}_{i-1},t) \qquad
\tilde{\rho}_i(x,0)=0. \label{rhodiffeq}
\end{align}

To show the existence of $\tilde{\rho}_i$ by induction, it suffices
to prove the following statement: For any $i \geq 1$, if
$|\tilde{\rho}_{i-1}| , |\nabla \tilde{\rho}_{i-1}| $ are bounded,
then $\tilde{\rho}_i$ exists and $|\tilde{\rho}_{i}|, |\nabla
\tilde{\rho}_{i}|$ are bounded. Assume  $|\tilde{\rho}_{i-1}| \leq
C_1, |\nabla \tilde{\rho}_{i-1}| \leq C_2$ are bounded on
$\mathbb{R}^{n+2} \times [0,T]$; then it follows from
Claim~\ref{claim:harmonic-map-flow-asymptotics} that $G(\tilde{\rho}_{i-1},{\bf x},t)$ is bounded
on $\mathbb{R}^{n+2} \times [0,T]$

\[
|G(\tilde{\rho}_{i-1},x,t)| \leq C_*(C_1,\tilde{h} ),
\]
and also because of Claim~\ref{claim:harmonic-map-flow-asymptotics} both $|\nabla B|$ and
$|\nabla \tilde h|$ are bounded on all of $\Ar^{n+2}\times [0,T]$,
it follows that $F(x,
\tilde{\rho}_{i-1},\nabla\tilde{\rho}_{i-1},t)$ is bounded:
\begin{align*}
& |F(x, \tilde{\rho}_{i-1},\nabla\tilde{\rho}_{i-1},t)|  \\
\leq & \left[(n-1)\sup |\nabla \tilde{h}| + \sup |\nabla B|
\right]C_2+C_2^2+ C_*(C_1,\tilde{h} ) = C_3
\end{align*}
Hence $\tilde{\rho}_i$ exists.

The bounds on $|\tilde{\rho}_{i}|$ and $|\nabla \tilde{\rho}_{i}|$
follow from the following estimates

\[
| \tilde{\rho}_i | \leq \int_0^t \int_{\mathbb{R}^{n+2}}
H(x,y,t-s)C_3 dyds \leq C_3t,
\]
and

\begin{align*}
& | \nabla \tilde{\rho}_i | = | \int_0^t \int_{\mathbb{R}^{n+2}}
[\nabla_x H(x,y,t-s)]
F(y,\tilde{\rho}_{i-1}, \nabla\tilde{\rho}_{i-1} ,t)dyds| \\
& \leq \int_0^t \int_{\mathbb{R}^{n+2}} | \nabla_x H(x,y,t-s)| C_3 dyds \\
& = \int_0^t \int_{\mathbb{R}^{n+2}} \frac{1}{(4 \pi
{(t-s)})^{(n+2)/2}}
e^{-\frac{|x-y|^2}{4{(t-s)}} } \frac{|x-y|}{2{(t-s)}}  C_3 dyds \\
& \leq \frac{(n+2)C_3}{\sqrt{\pi}} \int_0^t  \frac{1}{\sqrt{t-s}} ds
= \frac{2(n+2)C_3}{\sqrt{\pi}} \sqrt{t}.
\end{align*}

Assuming, as we shall, that $T\le   \min \{\frac{C_3}{C_1},\frac{\pi
C_2^2}{ 4(n+2)^2C_3^2} \}$, then for $0 \leq t \leq T$ we have for
all $i$,
\begin{align}
|\tilde{\rho}_{i}| \leq C_1 \qquad \text{and} \qquad  |\nabla
\tilde{\rho}_{i}| \leq C_2. \label{rho_i est}
\end{align}

We prove the convergence of $\tilde{\rho}_i$ to a solution of
(\ref{betterharMapeq}) via proving that it is a Cauchy sequence in
$C^1$-norm.  Note that $\tilde{\rho}_i- \tilde{\rho}_{i-1}$
satisfies

\begin{align}
& \frac{\partial (\tilde{\rho}_i- \tilde{\rho}_{i-1})}{\partial t} =
\Delta (\tilde{\rho}_i- \tilde{\rho}_{i-1}) + F(x,
\tilde{\rho}_{i-1},\nabla\tilde{\rho}_{i-1},t)-F(x,\tilde{\rho}_{i-2},
\nabla\tilde{\rho}_{i-2},t) \nonumber \\
& (\tilde{\rho}_i- \tilde{\rho}_{i-1})(x,0)=   0. \label{rhoiminus}
\end{align}
where
\begin{align*}
& F(x,
\tilde{\rho}_{i-1},\nabla\tilde{\rho}_{i-1},t)-F(x,\tilde{\rho
}_{i-2},\nabla\tilde{\rho}_{i-2},t) \\
= & [(n-1) \nabla \tilde{h}- \nabla B +\nabla (\tilde{\rho}_{i-1}+
\tilde{\rho}_{i-2})]
 \cdot \nabla(\tilde{\rho}_{i-1} - \tilde{\rho}_{i-2}) \\
& + G( \tilde{\rho}_{i-1},{\bf x},t)-G(\tilde{\rho}_{i-2},{\bf x},t)
\end{align*}

By lengthy but straightforward calculations one can verify the
Lipschitz property of $G(\tilde{\rho},{\bf x},t)$

\[
|G(\tilde{\rho}_{i-1},{\bf x},t)-G(\tilde{\rho}_{i-2},{\bf x},t)|
\leq C_\&(C_1,C _2, \tilde{f},\tilde{f}_0) \cdot |
\tilde{\rho}_{i-1}-\tilde{\rho}_{i-2}|.
\]
This and (\ref{rho_i est}) implies

\begin{align}
& |F(x,
\tilde{\rho}_{i-1},\nabla\tilde{\rho}_{i-1},t)-F(x,\tilde{\rho}_{i-2},
\nabla\tilde{\rho}_{i-2},t)|
\nonumber \\
\leq &  C_4 \cdot | \tilde{\rho}_{i-1}-\tilde{\rho}_{i-2}| +C_5
\cdot |\nabla \tilde{\rho}_{i-1}- \nabla \tilde{\rho}_{i-2} |
\label{Flip}
\end{align}
where $C_4 = C_\&(C_1,C_2, \tilde{f},\tilde{f}_0)$ and $C_5 = [(n-1)
\sup |\nabla \tilde{f}| +\sup |\nabla B|+2C_2]$.

Let

\begin{align*}
& A_i(t)= \sup_{0 \leq s \leq t, x \in \mathbb{R}^{n+2}}
|\tilde{\rho}_i-
\tilde{\rho}_{i-1}|(x,s) \\
& B_i(t)=\sup_{0 \leq s \leq t,x \in \mathbb{R}^{n+2}}
|\nabla(\tilde{\rho}_i-
 \tilde{\rho}_{i-1})|(x,s).
\end{align*}

From Equations~(\ref{rhoiminus}) and~(\ref{Flip}) we can estimate
$|\tilde{\rho}_i- \tilde{\rho}_{i-1}|$ and $ |\nabla(\tilde{\rho}_i-
\tilde{\rho}_{i-1})|$ in the same way as we estimate
$|\tilde{\rho}_{i}|$ and $|\nabla\tilde{\rho}_i|$ above; we conclude

\begin{align*}
& A_i(t)\leq [ C_4  A_{i-1}(t) +C_5 B_{i-1}(t)] \cdot t  \\
& B_i(t) \leq \frac{2(n+2)[C_4  A_{i-1}(t) +C_5
B_{i-1}(t)]}{\sqrt{\pi}} \cdot \sqrt{t}.
\end{align*}

Let $C_6 = \max \{C_4,C_5 \}$; then we get

\[
A_i(t)+B_i(t) \leq \left( C_6t+ \frac{2(n+2)C_6
\sqrt{t}}{\sqrt{\pi}} \right) \cdot \left(  A_{i-1}(t)+B_{i-1}(t)
\right).
\]

Now suppose that $T\le T_2$ where $T_2$ satisfies  $C_6T_2+
\frac{2(n+2)C_6 \sqrt{T_2}}{\sqrt{\pi}} =\frac{1}{2}$; then for all
$t\le T$ we have
\[
A_i(t)+B_i(t) \leq \frac{1}{2} \left(  A_{i-1}(t)+B_{i-1}(t)
\right).
\]
This proves that $\tilde{\rho}_{i}$ is a Cauchy sequence in
$C^1(\mathbb{R}^{n+2})$. Let $\lim_{i \rightarrow +\infty}
\tilde{\rho}_{i} =\tilde{\rho}_{\infty}$. Then $\nabla
\tilde{\rho}_{i} \rightarrow \nabla \tilde{\rho}_{\infty}$ and
$F(x,\tilde{\rho}_{i-1},\nabla \tilde{\rho}_{i-1},t) \rightarrow
F(x, \tilde{\rho}_{\infty}, \nabla \tilde{\rho}_{\infty},t)$
uniformly. Hence we get from (\ref{rhoIntg}),

\begin{align} \tilde{\rho}_\infty=\int_0^t \int_{\mathbb{R}^{n+2}} H(x,y ,t-s)F(y,\tilde{ \rho}_{\infty},
\nabla\tilde{\rho}_{\infty} ,t)dyds \label{rhoInfIntg}
\end{align}

The next argument is similar to the argument in \cite{LiTam}, p.21.
The function $\tilde{\rho}_{i}$ is a smooth solution of
(\ref{rhodiffeq}) with $\tilde{\rho}_i(x,0)=0$. Also, both
 $\tilde{\rho}_{i}$
and $F(x,\tilde{\rho}_{i_1},\nabla \tilde{\rho}_{i-1},t)$ are
uniformly bounded on $\mathbb{R}^{n+2} \times[0,T]$. Thus, by
Theorem 1.11 \cite{LSU}, p.211 and Theorem 12.1 \cite{LSU}, p.223,
for any compact $K \subset \mathbb{R}^{n+2}$ and any $0<t_*<T$,
there is $C_7$ and $\alpha \in(0,1)$ independent of $i$ such that

\[
| \nabla \tilde{\rho}_{i}(x,t) -\nabla \tilde{\rho}_{i}(y,s)| \leq
C_7 \cdot \left( |x-y|^\alpha+ |t-s|^{\alpha/2} \right)
\]
where $x,y \in K$ and $0 \leq t< s \leq t_*$.

Letting $i \rightarrow \infty$ we get

\begin{align}
| \nabla \tilde{\rho}_{\infty}(x,t) -\nabla
\tilde{\rho}_{\infty}(y,s)| \leq C_7 \cdot \left( |x-y|^\alpha+
|t-s|^{\alpha/2} \right). \label{rhoHolder}
\end{align}

Hence $\nabla \tilde{\rho}_{\infty}\in C^{\alpha,\alpha/2}$, i.e.,
it  is $\alpha$-H\"{o}lder continuous in space and
$\alpha/2$-H\"older continuous.

From (\ref{rhoInfIntg}) we conclude that $\tilde{\rho}_{\infty}$ is
a solution of (\ref{betterharMapeq}) on $\mathbb{R}^{n+2} \times
[0,T]$ with $\tilde{\rho}_{\infty}(x,0)=0$.





\subsection{The asymptotic behavior of the
solutions}\label{harmsec}

In the rest of this subsection we study the asymptotic behavior of
solution $\tilde{\rho}(x,t)$  as $x \rightarrow \infty$. First we
prove inductively that there is a constant $\lambda$ and $T_3$ such
that, provided that $T\le T_3$, for $x \in \mathbb{R}^{n+2}, t \in
[0,T]$, we have

\begin{align}
| \tilde{\rho}_i (x,t) | \leq  \frac{\lambda }{(1+|x|)^2} \qquad
\text{and} \qquad |\nabla \tilde{\rho}_i (x,t) | \leq
\frac{\lambda}{(1+|x|)^2} \label{rho_iDcay}
\end{align}

Clearly, since $\tilde \rho_0=0$, these estimates hold for $i=0$. It
follows from (\ref{rho_i est}) and Claim~\ref{claim:harmonic-map-flow-asymptotics} that there is
a constant $C_8$ independent of $i$ such that

\begin{align*}
& | G(\tilde{\rho}_i,{\bf x},t) | \leq \frac{C_8}{(1+|x|)^2} \\
& \left[ (n-1) |\nabla \tilde{h}|+ |\nabla B| \right](x,t) \leq C_8.
\end{align*}



Now we assume these estimates hold for $i$. Then for $0 \leq t \leq
T$ we have

\begin{align*}
| \tilde{\rho}_i (x,t) | & \leq  \int_0^t \int_{\mathbb{R}^{n+2}}
H(x,y,t-s) \left[ \frac{C_8 \lambda }{(1+|y|)^2} + \frac{\lambda^2
}{(1+|y|)^2}
+\frac{C_8}{(1+|y|)^2} \right ] dyds \\
& = \int_0^t \int_{\mathbb{R}^{n+2}}
\frac{1}{(4\pi(t-s))^{(n+2)/2}}e^{ -\frac{|x-y|^2}{4(t-s)}} \left[
\frac{C_8 \lambda +\lambda^2+C_8}{(1+|y|)^2}  \right ] dyds \\
& \leq (C_8 \lambda +\lambda^2+C_8) \cdot \frac{C(n)t}{(1+|x|)^2}.
\end{align*}

Also, we have

\begin{align*}
& |\nabla \tilde{\rho}_i (x,t) | \leq  \int_0^t
\int_{\mathbb{R}^{n+2}} |\nabla_x H(x,y,t-s)| \left[
\frac{C_8 \lambda + \lambda^2 + C_8}{(1+|y|)^2} \right ]  dyds \\
& = \int_0^t \int_{\mathbb{R}^{n+2}}
\frac{|x-y|}{2(t-s)}\frac{1}{(4\pi(t-s))^{(n+2)/2}}e^{-\frac{|x-y|^2}{4(t-s)}}
\left[
\frac{C_8 \lambda +\lambda^2+C_8}{(1+|y|)^2}  \right ] dyds \\
& \leq (C_8 \lambda +\lambda^2+C_8) \cdot
\frac{C(n)\sqrt{t}}{(1+|x|)^2}.
\end{align*}

If we choose $T_3$ such that

\[
(C_8 \lambda +\lambda^2+C_8) \cdot C(n)T_3 \leq \lambda \quad
\text{and} \quad (C_8 \lambda +\lambda^2+C_8) \cdot C(n) \sqrt{T_3}
\leq \lambda,
\]
then (\ref{rho_iDcay}) hold for all $i$. From the definition of
$\tilde{ \rho}_\infty$ we conclude

\begin{align}
| \tilde{\rho}_\infty (x,t) | \leq  \frac{\lambda }{(1+|x|)^2}
\qquad \qquad |\nabla \tilde{\rho}_\infty (x,t) | \leq
\frac{\lambda}{(1+|x|)^2} \label{rhodecay}
\end{align}


Recall that $\tilde\rho_\infty$ is a solution of the following
linear equation (in $\upsilon$):

\begin{align*}
& \frac{\partial \upsilon }{\partial t} = \Delta \upsilon + \nabla [
(n-1)\tilde{h} -B] \cdot \nabla \upsilon
+G( \tilde{\rho}_\infty,{\bf x},t) \\
& \upsilon(x,0)=0.
\end{align*}

From (\ref{rhoHolder}) and Claim~\ref{claim:harmonic-map-flow-asymptotics} we know that $\nabla [
(n-1)\tilde{h} -B+ \tilde{\rho}_\infty ]$ has
$C^{\alpha,\alpha/2}$-H\"older-norm bounded (this means
$\alpha$-H\"older norm in space and the $\alpha/2$-H\"older norm in
time). By some lengthy calculation we
 get

\[
| G(\tilde{\rho}_\infty,{\bf x},t)|_{C^{\alpha,\alpha/2}} \leq
\frac{C_9}{(1+|x|)^2}.
\]

By local Schauder estimates for parabolic equations we conclude

\[
|\tilde{\rho}_\infty |_{C^{2+\alpha,1+\alpha/2}} \leq
\frac{C_{10}}{(1+|x|)^2}.
\]

Using this estimate one can further show by calculation that

\begin{align*}
& |\nabla \nabla [ (n-1)\tilde{f} -B+
\tilde{\rho}_\infty ] |_{C^{\alpha,\alpha/2}} \leq C_{11} \\
& |\nabla G(\tilde{\rho}_\infty,{\bf x},t)|_{C^{\alpha,\alpha/2}}
\leq \frac{C_{12}}{(1+|x|)^2}.
\end{align*}

By local high order Schauder estimates for parabolic equations we
conclude

\[
|\nabla \tilde{\rho}_\infty |_{C^{2+\alpha,1+\alpha/2}} \leq
\frac{C_{13}}{(1+|x|)^2}.
 \]

We have proved the following:

\begin{proposition}[Harmonic map flow with quadratic decay at infinity]
\label{prop:harmonic-map-flow-existence-decay}
\label{harmapflow}
\uses{prop:harmonic-map-flow-short-time-existence,claim:harmonic-map-flow-asymptotics}
For a standard solution $(\mathbb{R}^n,g(t)),\ 0\le t<T$, and for
any $t_0\in [0,T)$ there is a rotationally symmetric solution
$\psi_t({\bf x)} = x e^{\tilde{\rho}({\bf x},t)}$ to the harmonic
map flow
\[
\frac{\partial \psi_t}{\partial t}=\Delta_{g(t),g(t_0)}\psi(t)
\qquad \psi(t_0)({\bf x})={\bf x},
\]
and $ |\nabla^i \tilde{\rho}|({\bf x},t)\leq \frac{C_{14}}{(1+|{\bf
x}|)^2}$ for $0 \leq i \leq 3$ defined on some non-degenerate
interval $[t_0,T']$.
\end{proposition}





\subsection{The uniqueness for the solutions of Ricci-DeTurck flow}

We prove the following general uniqueness result for Ricci-DeTurck
flow on open manifolds.

\begin{proposition}[Uniqueness for the Ricci-DeTurck flow on noncompact manifolds]
\label{prop:ricci-deturck-uniqueness-noncompact}
\label{RDuniq}
\uses{def:standard-ricci-flow}
Let $\hat g_1(  t)  $ and $\hat g_2(  t),\ 0 \leq t\leq T $, be two
bounded solutions of the Ricci-DeTurck flow on complete and
noncompact manifold $M^{n}$ with initial metric $g_1( t_0) =g_2(t_
0)  =g$. Suppose that for some $1<C<\infty$ we have
\begin{align*}
& C^{-1}  g    \leq \hat g_1( t) \leq C  g \\
& C^{-1}  g   \leq \hat g_2(  t) \leq C g .\\
\end{align*}
Suppose that in addition we have


\begin{align*}
& \left\Vert \hat g_1(t) \right\Vert _{C^{2}\left(M\right),g} \le C \\
& \left\Vert \hat g_2(t) \right\Vert _{C^{2}\left(M\right),g} \le
C.
\end{align*}




Lastly, suppose there is an exhausting sequence of compact, smooth
submanifolds of $\Omega_{k}\subset M$, i.e., $\Omega _{k}\subset\operatorname{int}\Omega_{k+1}$ and $\cup\Omega_{k}=M$ such that $\hat g_1\left(
t\right)$ and $\hat g_2(t)$ have the same sequential asymptotic
behavior at $\infty$ in the sense that for any $\epsilon>0,$ there
is a $k_0$ arbitrarily large with

\[
\left\vert \hat g_1(  t)  -\hat g_2(  t) \right\vert
_{C^{1}\left(\partial\Omega_{k_0}\right) ,g}\leq\epsilon,
\]
Then $ \hat g_1(  t)  =\hat g_2(  t)$.
\end{proposition}

\begin{proof}
Letting $\tilde \nabla$ be the covariant derivative determined by
$g$, then, using the Ricci-DeTurck flow (\ref{eqRcDeTurck}) for
$\hat g_1$ and $\hat g_2$, we can  make the following estimate for
an appropriate constant $C_{14}$ depending on $g$.
\begin{align*}
& \frac{\partial}{\partial t}\left\vert \hat g_1(  t) -\hat g_2( t)
\right\vert _{g}^{2} =2\left\langle \frac{\partial}{\partial
t}\left(  \hat g_1(  t) -\hat g_2( t)  \right) ,\hat g_1( t)
-\hat g_2(  t)  \right\rangle _{g}   \\
& \leq 2\left\langle\hat
g_1^{\alpha\beta}\tilde{\nabla}_{\alpha}\tilde{\nabla}_{\beta}\left(
\hat g_1(t)-\hat g_2(t)\right),  \left(  \hat g_1( t) -\hat g_2(t)
\right)\right\rangle_g
\\ &  +C_{14}\left\vert \hat g_1(  t)  -\hat g_2(
t) \right\vert _{g}^{2}  +C_{14}\left\vert \tilde\nabla\left( \hat
g_1( t) - \hat g_2( t)\right)  \right\vert _{g} \left\vert \hat g_1(
t)
-\hat g_2(  t)  \right\vert _{g}\\
& \leq \hat
g_1^{\alpha\beta}\tilde{\nabla}_{\alpha}\tilde{\nabla}_{\beta}\left(\left\vert
\hat g_1(  t)  -\hat g_2(  t)  \right\vert _{g}^{2}\right) -2\hat
g_1^{\alpha\beta}\left\langle\tilde{\nabla}_{\beta}\left( \hat
g_1(t)-\hat g_2(t)\right),\tilde{\nabla}_{\alpha}\left(  \hat g_1(
t)  -\hat g_2 (t)  \right)\right\rangle_g \\
&  +C_{14}\left\vert \hat g_1(t)-\hat g_2(t) \right\vert _{g}^{2}
+C_{14}\left\vert \tilde\nabla\left(  \hat g_1(t)-\hat g_2(t)
\right) \right\vert
_{g}\left\vert \hat g_1(t) - \hat g_2(t)\right\vert _{g}\\
&  \leq   \hat
g_1^{\alpha\beta}\tilde{\nabla}_{\alpha}\tilde{\nabla}_{\beta}\left(\left\vert
\hat g_1\left(  t\right)  -\hat g_2\left(  t\right)  \right\vert
_{g}^{2}\right)-2C^{-1}  \left\vert \tilde\nabla\left( \hat
g_1(t)-\hat g_2(t)\right)
\right\vert _{g}^{2}\\
&  +C_{14}\left\vert \hat g_1(t)-\hat g_2(t)\right\vert
_{g}^{2}+C^{-1} \left\vert \tilde\nabla\left( \hat g_1(t)-\hat
g_2(t) \right) \right\vert _{g}^{2}+\frac{C_{14}^{2}}{4C^{-1}
}\left\vert \hat g_1(t)-\hat g_2(t) \right\vert _{g}^{2},
\end{align*}
where the last inequality comes from completing the square to
replace the last term in the previous expression. Thus, we have
proved

\begin{equation}
\frac{\partial}{\partial t}\left\vert \hat g_1\left(  t\right) -\hat
g_2\left( t\right) \right\vert _{g}^{2}\leq 2 \hat
g_1^{\alpha\beta}\tilde{\nabla}_{\alpha}
\tilde{\nabla}_{\beta}\left\vert \hat g_1\left(  t\right)  -\hat
g_2\left( t\right) \right\vert _{g}^{2}+C_{15}\left\vert \hat
g_1\left( t\right)  -\hat g_2\left( t\right)  \right\vert _{g}^{2}
\label{eqestforunique}
\end{equation}
pointwise on $\Omega_{k}$ with $C_{15}$ a constant that depends only
on $n$, $C$ and $g$.

Suppose that $\hat g_1\left(  t\right)  \neq  \hat g_2 \left(
t\right) $ for some $t$. Then there is a point $x_{0}$ such that
$\left\vert \hat g_1(x_0, t) - \hat g_2(x_0, t) \right\vert _{g}^{2}
>\epsilon_{0}$ for  some $\epsilon_{0}>0$.

We choose a $k_0$ sufficiently large that $x_{0}\in \Omega_{k_0}$
and for all $t'\in [t_0,T]$ we have

\begin{equation}
\sup_{x\in\partial\Omega_{b}}\left\vert \hat g_1(x,  t') -\hat
g_2(x, t') \right\vert _{g}^{2} \leq \epsilon \label{eq est for uniq
bddy}
\end{equation}
where $\epsilon >0$ is a constant to be chosen later.

Recall we have the initial condition $\left\vert \hat g_1( 0) -\hat
g_2(0)  \right\vert _{g}^{2}=0$. Using Equation~(\ref{eqestforunique}) and applying the maximum principle to $\left\vert \hat g_1(
t) -\hat g_2( t) \right\vert _{g}^{2}$ ) on the domain
$\Omega_{k_0}$, we get
\[
e^{-C_{15}t}\left\vert \hat g_1 (  t)  -\hat g_2 ( t) \right\vert
_{g}^{2}(  x) \leq\epsilon.\text{ for all }x\in\Omega _{k_0}.
\]

This is a contradiction if we choose $\epsilon \leq
\epsilon_{0}e^{-C_{15}T}$. This contradiction establishes the
proposition.
\end{proof}

 Let $g_1(t),\ 0\le t<T_1$, and $g_2(t),\ 0\le t<T_2$, be standard solutions that
 agree on the interval $[0,t_0]$ for some $t_0\ge 0$.
By Proposition~\ref{prop:harmonic-map-flow-existence-decay} there are $\psi_1(t)$  and
$\psi_2(t)$  which are solutions of the harmonic map flow defined
for $t_0 \leq t \leq T$ for some $T>t_0$ for the Ricci flows
$g_1(t)$ and $g_2(t)$. Let $\hat{g}_1(t)= (\psi^{-1}(t))^*g_1(t)$
and $\hat{g}_2(t) = (\psi^{-1}(t))^*g_2(t)$. Then $\hat{g}_1(t)$ and
$\hat{g}_2(t)$ are two solutions of the Ricci-DeTurck flow with
$\hat{g}_1(t_0)=\hat{g}_2(t_0)$. Choose $T' \in (t_0,T]$ such that
$\hat{g}_1(t)$ and $\hat{g}_2(t)$ are $\delta$-close to
$\hat{g}_1(t_0)$ as required in Proposition~\ref{prop:ricci-deturck-uniqueness-noncompact}. It follows
from Lemma~\ref{prop:standard-flow-asymptotic-cylinder} and the decay estimate in
Proposition~\ref{prop:harmonic-map-flow-existence-decay} that $\hat{g}_1(t)$ and $\hat{g}_2(t)$
are bounded solutions and that they have same sequential asymptotic
behavior at infinity. We can apply Proposition~\ref{prop:ricci-deturck-uniqueness-noncompact} to
conclude $\hat{g}_1(t)=\hat{g}_2(t)$ on $t_0 \leq t \leq T'$.  We
have proved:

\begin{corollary}[Agreement of Ricci-DeTurck flows for standard solutions]
\label{cor:standard-solution-deturck-agreement}
\label{duniq}
\uses{prop:harmonic-map-flow-existence-decay,prop:ricci-deturck-uniqueness-noncompact}
Let $g_1(t)$ and $g_2(t)$  be standard solutions. Suppose that
$g_1(t)=g_2(t)$ for all $t\in [0,t_0]$ for some $t_0\ge 0$. The
Ricci-DeTurck solutions $\hat{g}_1(t)$ and $\hat{g}_2(t)$
constructed from standard solutions $g_1(t)$ and $g_2(t)$ with
$g_1(t_0)=g_2(t_0)$ exist and satisfy $\hat{g}_1(t)=\hat{g}_1(t)$
for $t \in [t_0,T']$ for some $T'>t_0$.
\end{corollary}




\section{Completion of the proof of uniqueness}

Now we are ready to prove the uniqueness of the standard solution.
Let $g_1(t),\ 0\le t<T_1$, and $g_2(t),\ 0\le t<T_2$, be a standard
solutions. Consider the maximal interval $I$ (closed or half-open)
containing $0$ on which $g_1$ and $g_2$ agree.

{\bf Case 1: $T_1< T_2$ and $I=[0,T_1)$}

In this case since $g_1(t)=g_2(t)$ for all $t<T_1$ and $g_2(t)$
extends smoothly past time $T_1$, we see that the curvature of
$g_1(t)$ is bounded as $t$ tends to $T_1$. Hence, $g_1(t)$ extends
past time $T_1$, contradicting the fact that it is a maximal flow.

{\bf Case 2: $T_2<T_1$ and $I=[0,T_2)$}

The argument in this case is the same as the previous one with the
roles of $g_1(t)$ and $g_2(t)$ reversed.

There is one more case to rule out.


{\bf Case 3: $I$ is a closed interval $I=[0,t_0]$.}

In this case, of course, $t_0<\min(T_1,T_2)$. Hence we apply
Proposition~\ref{prop:harmonic-map-flow-existence-decay} to construction solutions $\psi_1$ and
$\psi_2$ to the harmonic map flow for $g_1(t)$ and $g_2(t)$ with
$\psi_1$ and $\psi_2$ being the identity at time $t_0$. These
solutions will be defined on an interval of the form $[t_0,T]$ for
some $T>t_0$. Using these harmonic map flows we construct solutions
$\hat g_1(t)$ and $\hat g_2(t)$ to the Ricci-DeTurck flow defined on
the interval $[t_0,T]$. According to Corollary~\ref{cor:standard-solution-deturck-agreement}, there is
a uniqueness theorem for these Ricci-DeTurck flows, which implies
that $\hat g_1(t)=\hat g_2(t)$ for all $t\in [t_0,T']$ for some
$T'>t_0$. Invoking Corollary~\ref{cor:ricci-deturck-uniqueness-implies-ricci-uniqueness} we conclude that
$g_1(t)=g_2(t)$ for all $t\in [0,T']$, contradicting the maximality
of the interval $I$.

If none of these three cases can occur, then the only remaining
possibility is that $T_1=T_2$ and $I=[0,T_1)$, i.e., the flows are
the same. This then completes the proof of the uniqueness of the
standard flow.


\subsection{$T=1$ and existence of canonical neighborhoods}

At this point we have established all the properties claimed in
Theorem~\ref{thm:standard-solution-existence-uniqueness} for the standard flow except for the fact that
$T$, the endpoint of the time-interval of definition, is equal to
$1$. We have shown that $T\le 1$. In order to establish the opposite
inequality, we must show the existence of canonical neighborhoods
for the standard solution.



Here is the result about the existence of canonical neighborhoods
for the standard solution.

\begin{theorem}[Canonical neighborhoods in the standard solution]
\label{thm:standard-solution-canonical-neighborhoods}
\label{R_0cannbhd}
\uses{thm:standard-solution-existence-uniqueness,prop:standard-flow-kappa-noncollapsed}
Fix $0<\epsilon<1$.   Then there is $r>0$ such that for any point
$(x_0,t_0)$ in the standard flow with $R(x_0,t_0)\ge r^{-2}$ the
following hold.
\begin{enumerate}
\item[(1)] $t_0>r^2$.
\item[(2)] $(x_0,t_0)$ has a strong canonical
$(C(\epsilon),\epsilon)$-neighborhood. If this canonical neighborhood is a strong
$\epsilon$-neck centered at $(x_0,t_0)$, then the strong neck
extends to an evolving neck defined for backward rescaled time
$(1+\epsilon)$.
\end{enumerate}
\end{theorem}


\begin{proof}
\uses{prop:standard-flow-kappa-noncollapsed,kaplimit,smlmtflow,limitcannbhd}
 Take an increasing sequence of times $t'_n$ converging to $T$.
Since the curvature of $(\Ar^3,g(t))$ is locally bounded in time,
for each $n$, there is a bound on the scalar curvature on
$\Ar^3\times [0,t'_n]$. Hence, there is a finite upper bound $R_n$
on $R(x,t)$ for all points $(x,t)$ with $t\le t'_n$ for which the
conclusion of the theorem does not hold. (There clearly are such
points since the conclusion of the theorem fails for all $(x,0)$.)
Pick $(x_n,t_n)$ with $t_n\le t'_n$, with $R(x_n,t_n)\ge R_n/2$ and
such that the conclusion of the theorem does not hold for
$(x_n,t_n)$. To prove the theorem we must show that $\overline\lim_{n\rightarrow\infty}R(x_n,t_n)<\infty$. Suppose the contrary.
By passing to a subsequence we can suppose that $\lim_{n\rightarrow \infty}R(x_n,t_n)=\infty$. We set
$Q_n=R(x_n,t_n)$. We claim that all the hypotheses of
Theorem~\ref{kaplimit} apply to the sequence
$(\Ar^3,g(t),(x_n,t_n))$. First, we show that all the hypotheses of
Theorem~\ref{smlmtflow} (except the last) hold. Since $(\Ar^3,g(t))$
has non-negative curvature all these flows have curvature pinched
toward positive. By Theorem~\ref{prop:standard-flow-kappa-noncollapsed} there are $r>0$ and
$\kappa>0$ so that all these flows are $\kappa$-non-collapsed on
scales $\le r$. By construction if $t\le t_n$ and $R(y,t)>2Q_n\ge
R_n$ then  the point $(y,t)$ has a strong canonical
$(C(\epsilon),\epsilon)$-neighborhood. We are assuming that
$Q_n\rightarrow \infty$ as $n\rightarrow \infty$ in order to achieve
the contradiction. Since all time-slices are complete, all balls of
finite radius have compact closure.

Lastly, we need to show that the extra hypothesis of
Theorem~\ref{kaplimit} (which includes the last hypothesis of
Theorem~\ref{smlmtflow}) is satisfied. This is clear since
$t_n\rightarrow T$ as $n\rightarrow \infty$ and $Q_n\rightarrow
\infty$ as $n\rightarrow \infty$. Applying Theorem~\ref{kaplimit} we
conclude that after passing to a subsequence there is a limiting
flow which is a $\kappa$-solution. Clearly, this and
Corollary~\ref{limitcannbhd} imply that for all sufficiently large
$n$ (in the subsequence)  the neighborhood as required by the
theorem exists. This contradicts our assumption that none of the
points $(x_n,t_n)$ have these neighborhoods. This contradiction
proves the result.
\end{proof}

\subsection{Completion of the proof of
Theorem~{\ref{thm:standard-solution-existence-uniqueness}}}

The next proposition establishes the last of the conditions claimed
in Theorem~\ref{thm:standard-solution-existence-uniqueness}.


\begin{theorem}[The standard flow exists exactly on $[0,1)$]
\label{thm:standard-solution-time-interval-one}
\uses{thm:standard-solution-canonical-neighborhoods,prop:standard-flow-asymptotic-cylinder}
For the standard flow $T=1$.
\end{theorem}

\begin{proof}
\uses{cor:standard-flow-time-bound-one,thm:standard-solution-canonical-neighborhoods,prop:standard-flow-asymptotic-cylinder,shiw/deriv,kaplimit,asympvol,thm:bishop-gromov,lem:distance-nonincreasing}
We have already seen in Corollary~\ref{cor:standard-flow-time-bound-one} that $T\le 1$. Suppose now
that $T<1$. Take $T_0<T$ sufficiently close to $T$. Then according
to Proposition~\ref{prop:standard-flow-asymptotic-cylinder} there is a compact subset $X\subset
\Ar^3$ such that restriction of the flow to $(\Ar^3\setminus
X)\times [0,T_0]$ is $\epsilon$-close to the standard evolving flow
on $S^2\times (0,\infty),(1-t)h_0\times ds^2$, where $h_0$ is the
round metric of scalar curvature $1$ on $S^2$. In particular,
$R(x,T_0)\le (1+\epsilon)(1-T_0)^{-1}$ for all $x\in \Ar^3\setminus
X$. Because of Theorem~\ref{thm:standard-solution-canonical-neighborhoods} and the definition of
$(C(\epsilon),\epsilon)$-canonical neighborhoods, it follows that at
any point $(x,t)$ with $R(x,t)\ge r^{-2}$, where $r>0$ is the
constant given in Theorem~\ref{thm:standard-solution-canonical-neighborhoods}, we have $\partial
R/\partial t(x,t)\le C(\epsilon)R^2(x,t)$. Thus, provided that
$T-T_0$ is sufficiently small, there is a uniform bound to $R(x,t)$
for all $x\in \Ar^3\setminus X$ and all $t\in [T_0,T)$. Using
Theorem~\ref{shiw/deriv} and the fact that the standard flow is
$\kappa$-non-collapsed implies that the restrictions of the metrics
$g(t)$ to $\Ar^3\setminus X$ converge smoothly to a limiting
Riemannian metric $g(T)$ on $\Ar^3\setminus X$. Fix a non-empty open
subset $\Omega\subset \Ar^3\setminus X$ with compact closure. For
each $t\in [0,T)$ let $V(t)$ be the volume of
$(\Omega,g(t)|_{\Omega})$. Of course, $\lim_{t\rightarrow
T}V(t)=\Vol_{g(T)}(\Omega)>0$.




Since the metric $g(T)$ exists in a neighborhood of infinity and has
bounded curvature there, if the limit metric $g(T)$ exists on all of
$\Ar^3$, then we can extend the flow keeping the curvature bounded.
This contradicts the maximality of our flow subject to the condition
that the curvature be locally bounded in time. Consequently, there
is a point $x\in \Ar^3$ for which the limit metric $g(T)$ does not
exist. This means that $\overline\lim_{t\rightarrow
T}R(x,t)=\infty$. That is to say,  there is a sequence of
$t_n\rightarrow T$ such that setting $Q_n=R(x,t_n)$, we have
$Q_n\rightarrow \infty$ as $n$ tends to infinity. By
Theorem~\ref{thm:standard-solution-canonical-neighborhoods} the second hypothesis in the statement of
Theorem~\ref{smlmtflow} holds for the sequence
$(\Ar^3,g(t),(x,t_n))$. All the other hypotheses of this theorem as
well as the extra hypothesis in Theorem~\ref{kaplimit} obviously
hold for this sequence. Thus, according to Theorem~\ref{kaplimit}
the based flows $(\Ar^3,Q_ng(Q_n^{-1}t'+t_n),(x,0))$ converge
smoothly to a $\kappa$-solution. Since the asymptotic volume of any
$\kappa$-solution is zero (see Theorem~\ref{asympvol}), we see that
for all $n$ sufficiently large, the following holds:

\begin{claim}[Volume ratio of the rescaled limit is small]
\label{claim:kappa-solution-volume-ratio-small}
\uses{asympvol}
For any $\epsilon>0$, there is $A<\infty$ such that for all $n$
sufficiently large we have
$$\Vol(B_{Q_ng}(x,t_n,A))<\epsilon A^3.$$
\end{claim}


Rescaling, we see that for all $n$ sufficiently large we have
$$\Vol\,B_g(x,t_n,A/\sqrt{Q_n})< \epsilon(A/\sqrt{Q_n})^3.$$
Since the curvature of $g(t_n)$ is non-negative and since the $Q_n$
tend to $\infty$, it follows from the Bishop-Gromov Inequality
(Proposition~\ref{thm:bishop-gromov}) that for any $0<A<\infty$ and any
$\epsilon>0$, for all $n$ sufficiently large we have
$$\Vol\,B_g(x,t_n,A)<\epsilon A^3.$$

On the other hand, since $\Omega$ has compact closure, there is an
$A_1<\infty$ with $\Omega\subset B(x,0,A_1)$. Since the curvature of
$g(t)$ is non-negative for all $t\in [0,T)$, it follows from
Lemma~\ref{lem:distance-nonincreasing} that the distance is a non-increasing
function of $t$, so that for all $t\in [0,T)$ we have $\Omega\subset
B(x,t,A_1)$. Applying the above, for any $\epsilon>0$ for all $n$
sufficiently large we have
$$\Vol\,(\Omega,g(t_n))\le \Vol_g\,B(x,t_n,A_1)<\epsilon
(A_1)^3.$$ But this contradicts the fact that
$$\lim_{n\rightarrow
\infty}\Vol\,\Omega,g(t_n)=\Vol\,(\Omega,g(T))>0.$$ This
contradiction proves that $T=1$.
\end{proof}


This completes the proof of Theorem~\ref{thm:standard-solution-existence-uniqueness}.

\section{Some corollaries}

Now let us derive extra properties of the standard solution that
will be important in our applications.

\begin{proposition}[Scalar curvature lower bound $R\ge c/(1-t)$]
\label{prop:standard-solution-scalar-curvature-lower-bound}
\label{Restim}
\uses{thm:standard-solution-canonical-neighborhoods}
There is a constant $c>0$ such that for all $(p,t)$ in the standard
solution we have
$$R(p,t)\ge \frac{c}{1-t}.$$
\end{proposition}

\begin{proof}
\uses{thm:standard-solution-canonical-neighborhoods,lem:distance-nonincreasing,controlnbhd,canonvary,narrows,omega,kaplimit,asympvol,thm:bishop-gromov}
First, let us show that there is not a limiting metric $g(1)$
defined on all of $\Ar^3$. This does not immediately contradict the
maximality of the flow because we are assuming only that the flow is
maximal subject to having curvature locally bounded in time. Assume
that a limiting metric $(\Ar^3,g(1))$ exists.
 First,
notice that from the canonical neighborhood assumption and
Lemma~\ref{controlnbhd} we see that the curvature of $g(T)$ must be
unbounded at spatial infinity. On the other hand, by
Proposition~\ref{canonvary} every point of $(\Ar^3,g(1))$ of
curvature greater than $R_0$ has a $(2C,2\epsilon)$-canonical
neighborhood. Hence, since $(\Ar^3,g(1))$ has unbounded curvature,
it then has $2\epsilon$-necks of arbitrarily small scale. This
contradicts Proposition~\ref{narrows}. (One can also rule this
possibility out by direct computation using the spherical symmetry
of the metric.) This means that there is no limiting metric $g(1)$.



The next step is to see that for any $p\in \Ar^3$ we have $\lim_{t\rightarrow 1}R(p,t)=\infty$. Let $\Omega\subset \Ar^3$ be
the subset of $x\in \Ar^3$ for which $\liminf_{t\rightarrow
1}R(x,t)<\infty$. We suppose that $\Omega\not=\emptyset$. According
to Theorem~\ref{omega} the subset $\Omega$ is open and the metrics
$g(t)|_\Omega$ converge smoothly to a limiting metric
$g(1)|_\Omega$. On the other hand, we have just seen that there is
not a limit metric $g(1)$ defined everywhere. This means that there
is $p\in \Ar^3$ with  $\lim_{t\rightarrow 1}R(p,t)=\infty$.
Take a sequence $t_n$ converging to $1$ and set $Q_n=R(p,t_n)$. By
Theorem~\ref{kaplimit} we see that, possibly after passing to a
subsequence, the based flows $(\Ar^3,Q_ng(t'-t_n),(p,0))$ converge
to a $\kappa$-solution. Then by Proposition~\ref{asympvol} for any
$\epsilon>0$ there is $A<\infty$ such that $\Vol\,B_{Q_ng}(p,t_n,A)<\epsilon A^3$, and hence after rescaling we
have $\Vol\,B_g(p,t_n,A/\sqrt{Q_n})<\epsilon (A/\sqrt{Q_n})^3$.
By the Bishop-Gromov inequality (Proposition~\ref{thm:bishop-gromov}) it
follows that for any $0<A<\infty$, any $\epsilon>0$ and for all $n$
sufficiently large, we have $\Vol\,B_g(p,t_n,A)<\epsilon A^3$.
Take a non-empty subset $\Omega'\subset \Omega$ with compact
closure. Of course, $\Vol\,(\Omega',g(t))$ converges to $\Vol\,(\Omega',g(T))>0$ as $t\rightarrow T$. Then there is
$A<\infty$ such that for each $n$, the subset $\Omega'$ is contained
in the ball $B(p_0,t_n,A)$. This is a contradiction since it implies
that for any $\epsilon>0$ for all $n$ sufficiently large we have
$\Vol\,(\Omega',g(t))<\epsilon A^3$. This completes the proof
that for every $p\in \Ar^3$ we have $\lim_{t\rightarrow
1}R(p,t)=\infty$.

Fix $\epsilon>0$ sufficiently small and set $C=C(\epsilon)$. Then
for every $(p,t)$ with $R(p,t)\ge r^{-2}$ we have
$$\left|\frac{dR}{dt}(p,t)\right|\le CR^2(p,t).$$
 Fix $t_0=1-1/2r^2C$. Since the flow has curvature locally bounded in
 time, there is $2C\le C'<\infty$ such that $R(p,t_0)\le 1/(C'(1-t_0)$ for
 all $p\in \Ar ^3$.
 Since $R(p,t_0)=1/C'(1-t_0)$,
for all $t\in [t_0,1)$ we have
$$R(p,t)<\max\left(\left[(C'-C)(1-t_0)\right]^{-1},\left[r^{-2}-C(1-t_0)\right]^{-1}\right).$$

This means that $R(p,t)$ is uniformly bounded as $t\rightarrow 1$,
contradicting what we just established. This shows that for $t\ge
1-1/2r^2C$ the result holds. For $t\le 1-1/2r^2C$ there is a
positive lower bound on the scalar curvature, and hence the result
is immediate for these $t$ as well.
\end{proof}

\begin{theorem}[Trichotomy of canonical neighborhoods in the standard solution]
\label{thm:standard-solution-trichotomy}
\label{stdsolncannbhd}
\uses{thm:standard-solution-canonical-neighborhoods,prop:standard-solution-scalar-curvature-lower-bound}
For any $\epsilon>0$ there is $C'(\epsilon)<\infty$ such that for
any point $x$ in the standard solution one of the following holds.
\begin{enumerate}
\item[(1)] $(x,t)$ is contained in the core of a  $(C'(\epsilon),\epsilon)$-cap.
\item[(2)] $(x,t)$ is the center of an evolving $\epsilon$-neck $N$ whose initial time-slice is
$t=0$ and whose initial time-slice is disjoint from the surgery cap.
\item[(3)] $(x,t)$ is the center of an evolving $\epsilon$-neck defined for rescaled time
$1+\epsilon$.
\end{enumerate}
\end{theorem}



\begin{remark}[Evolving necks of rescaled time greater than one]
\label{rem:evolving-neck-rescaled-time}
\uses{def:evolving-epsilon-neck}
At first glance it may seem impossible for a point $(x,t)$ in the
standard solution to be the center of an evolving $\epsilon$-neck
defined for rescaled time $1+\epsilon$ since the standard solution
itself is only defined for time $1$. But this is indeed possible.
The reason is because the scale referred to for an evolving neck
centered at $(x,t)$ is $R(x,t)^{-1/2}$. As $t$ approaches one,
$R(x,t)$ goes to infinity, so that rescaled time $1$ at $(x,t)$ is
an arbitrarily small time interval measured in the scale of the
standard solution.
\end{remark}



\begin{proof}
\uses{thm:standard-solution-canonical-neighborhoods,prop:standard-solution-scalar-curvature-lower-bound,prop:standard-flow-asymptotic-cylinder}
By Theorem~\ref{thm:standard-solution-canonical-neighborhoods}, there is $r_0$ such that if $R(x,t)\ge
r_0^{-2}$, then $(x,t)$ has a  $(C,\epsilon)$-canonical neighborhood
and if this canonical neighborhood is a strong $\epsilon$-neck
centered at $x$, then that neck extends to an evolving neck defined
for rescaled time $(1+\epsilon)$. By Proposition~\ref{prop:standard-solution-scalar-curvature-lower-bound}, there
is $\theta<1$ such that if $R(x,t)\le r_0^{-2}$ then $t\le \theta$.
By Proposition~\ref{prop:standard-flow-asymptotic-cylinder},  there is a compact subset $X\subset \Ar
^3$ such that if $t\le \theta$ and $x\notin X$, then there is an
evolving $\epsilon$-neck centered at $x$ whose initial time is zero
and whose initial time-slice is at distance at least one from the
surgery cap. Lastly, by compactness there is $C'<\infty$ such that
every $(x,t)$ for $x\in X$ and every $t\le \theta$ is contained in
the core of a $(C',\epsilon)$-cap.\end{proof}

\begin{corollary}[Points near the tip lie in a canonical cap]
\label{cor:standard-solution-cap-near-tip}
\label{stdsolncap}
\uses{thm:standard-solution-trichotomy}
Fix $\epsilon>0$. Suppose that $(q,t)$ is a point in the standard
solution with $t\le R(q,t)^{-1}(1+\epsilon))$ and with $(q,0)\in
B(p_0,0,(\epsilon^{-1}/2)+A_0+5)$. Then $(q,t)$ is contained in an
$(C'(\epsilon),\epsilon)$-cap.
\end{corollary}

\begin{remark}[The tip $p_0$ and the constant $A_0$]
\label{rem:tip-and-a0-definitions}
\uses{lem:standard-initial-metric-asymptotic-bounds}
Recall that $p_0$ is the origin in $\Ar^3$ and hence is the tip of
the surgery cap. Also, $A_0$ is defined in Lemma~\ref{lem:standard-initial-metric-asymptotic-bounds}.
\end{remark}


\begin{corollary}[Trichotomy for flows converging to the standard solution]
\label{cor:generalized-flow-limit-trichotomy}
\label{stdsolnlimit}
\uses{thm:standard-solution-trichotomy}
For any $\epsilon>0$ let $C'=C'(\epsilon)$ be as in
Theorem~\ref{thm:standard-solution-trichotomy}. Suppose that we have a sequence of
generalized Ricci flows $({\mathcal M}_n,G_n)$, points $x_n\in
{\mathcal M}_n$ with ${\bf t}(x_n)=0$, neighborhoods $U_n$ of $x_n$
in the zero time-slice of ${\mathcal M}_n$, and a constant
$0<\theta<1$. Suppose that there are embeddings $\rho_n\colon
U_n\times [0,\theta)\to {\mathcal M}_n$ compatible with time and the
vector field so that the Ricci flows $\rho_n^*G_n$ on $U_n$ based at
$x_n$ converge geometrically to the restriction of the standard
solution to $[0,\theta)$. Then for all $n$ sufficiently large, and
any point $y_n$ in the image of $\rho_n$ one of the following holds:
\begin{enumerate}
\item[(1)] $y_n$ is contained in the core of a
$(C'(\epsilon),\epsilon)$-cap
\item[(2)] $y_n$ is the center of a strong $\epsilon$-neck
\item[(3)] $y_n$ is the center of an evolving $\epsilon$-neck whose
initial time-slice is at time $0$.
\end{enumerate}
\end{corollary}

\begin{proof}
\uses{thm:standard-solution-trichotomy,canonvary}
This follows immediately from Theorem~\ref{thm:standard-solution-trichotomy} and
Proposition~\ref{canonvary}.
\end{proof}



There is one property that we shall use later in proving the
finite-time extinction of Ricci flows with surgery for manifolds
with finite fundamental group (among others). This is a distance
decreasing property which we record here.

Notice that for the standard initial metric constructed in
Lemma~\ref{lem:existence-standard-initial-metric} we have the following:

\begin{lemma}[Distance-decreasing property of the standard initial metric]
\label{lem:standard-initial-metric-distance-decreasing}
\label{decrease}
\uses{def:standard-initial-metric}
Let $S^2$ be the unit sphere in $T_0\Ar^3$. Equip it with the metric
$h_0$ that is twice the usual metric (so that the scalar curvature
of $h_0$ is $1$). We define a map $\rho\colon S^2\times
[0,\infty)\to \Ar^3$ by sending the point $(x,s)$ to the point at
distance $s$ from the origin in the radial direction from $0$ given
by $x$ (all this measured in the metric $g_0$). Then $\rho^*g_0\le
h_0\times ds^2$.
\end{lemma}


\begin{proof}
Clearly, the metric $\rho^*g_0$ is rotationally symmetric and its
component in the $s$-direction is $ds^2$. On the other hand, since
each cross section $\{s\}\times S^2$ maps conformally onto a sphere
of radius $\le\sqrt{2}$ the result follows.
\end{proof}
